\subsection{Convolution and linear representations}

PROPOSITION 11. --- Let G be a locally compact group, E a quasi-complete locally convex space, U a continuous representation of G in E.

\begin{enumerate}
\def\labelenumi{(\roman{enumi})}
\item
  If \(\lambda \in \mathcal{C}'(\mathbf{G})\) , \(\mu \in \mathcal{C}'(\mathbf{G})\) , then \(U(\lambda * \mu) = U(\lambda)U(\mu)\) .
\item
  Suppose that \(\mathbf{E}\) is a Banach space, and let \(\rho(s) = \|U(s)\|\) for \(s\in\mathbf{G}\) . If \(\lambda\in\mathcal{M}^{\rho}(\mathbf{G})\) , \(\mu\in\mathcal{M}^{\rho}(\mathbf{G})\) , then \(U(\lambda*\mu)=U(\lambda)U(\mu)\) .
\end{enumerate}

Let \(\lambda, \mu\) be in \(\mathcal{C}'(\mathrm{G})\) . For any \(x \in \mathrm{E}\) one has, by applying notably Props. 1 and 4 of Ch. VI, §1, No.~1,

\[
\begin{array}{r l} & U (\lambda * \mu) x = \int_ {\mathbf {G}} U (s) x d (\lambda * \mu) (s) \\ & \quad = \int_ {\mathbf {G} \times \mathbf {G}} U (s t) x d \lambda (s) d \mu (t) = \int_ {\mathbf {G} \times \mathbf {G}} U (s) U (t) x d \lambda (s) d \mu (t) \\ & \quad = \int_ {\mathbf {G}} U (\lambda) U (t) x d \mu (t) = U (\lambda) \int_ {\mathbf {G}} U (t) x d \mu (t) = U (\lambda) U (\mu) x, \end{array}
\]

whence (i). An analogous argument may be applied in case (ii).

With G still a locally compact group, let us assume that G operates continuously on the left in a locally compact space X. This defines (§2, No.~4) a continuous linear representation \(\pmb{\gamma}\) of G in \(\mathcal{M}(\mathrm{X})\) (equipped with the topology of compact convergence in \(\mathcal{K}(\mathrm{X})\) ).

PROPOSITION 12. --- If \(\lambda \in \mathcal{C}'(\mathbf{G})\) and \(\mu \in \mathcal{M}(\mathbf{X})\) , then

\[
\gamma (\lambda) \mu = \lambda * \mu .
\]

By Prop. 7 of §1, No.~5,

\[
\lambda * \mu = \int_ {G} (\varepsilon_ {s} * \mu) d \lambda (s).
\]

Now, \(\varepsilon_{s}*\mu = \pmb {\gamma}(s)\mu\) (No.~2, formula (5)) and

\[
\int_ {\mathsf {G}} \left(\pmb {\gamma} (s) \mu\right) d \lambda (s) = \pmb {\gamma} (\lambda) \mu
\]

by the definition of \(\pmb{\gamma}(\lambda)\) .

COROLLARY. --- The mapping \((\lambda,\mu)\mapsto\lambda*\mu\) of \(\mathcal{C}^{\prime}(\mathrm{G})\times\mathcal{M}(\mathrm{X})\) into \(\mathcal{M}(\mathrm{X})\) is hypocontinuous relative to the equicontinuous subsets of \(\mathcal{C}^{\prime}(\mathrm{G})\) and the compact subsets of \(\mathcal{M}(\mathrm{X})\) ( \(\mathcal{C}^{\prime}(\mathrm{G})\) and \(\mathcal{M}(\mathrm{X})\) being equipped with the topology of compact convergence in \(\mathcal{C}(\mathrm{G})\) and \(\mathcal{H}(\mathrm{X})\) , respectively).

For, \(\mathcal{M}(\mathrm{X})\) , equipped with the topology of compact convergence in \(\mathcal{K}(\mathrm{X})\) , is quasi-complete. Therefore the mapping \((\lambda, \mu) \mapsto \pmb{\gamma}(\lambda)\mu\) of \(\mathcal{C}'(\mathrm{G}) \times \mathcal{M}(\mathrm{X})\) into \(\mathcal{M}(\mathrm{X})\) is hypocontinuous relative to the equicontinuous subsets of \(\mathcal{C}'(\mathrm{G})\) and the compact subsets of \(\mathcal{M}(\mathrm{X})\) (§2, No.~6). It then suffices to apply Prop. 12.

Remarks. --- 1) Let \(\lambda_0 \in \mathcal{C}'(\mathbf{G})\) . The mapping \(\mu \mapsto \lambda_0 * \mu\) of \(\mathcal{M}(\mathbf{X})\) into \(\mathcal{M}(\mathbf{X})\) is vaguely continuous. For, let \(f \in \mathcal{K}(\mathbf{X})\) . One has

\[
\langle \lambda_ {0} * \mu , f \rangle = \int f (s x) d \lambda_ {0} (s) d \mu (x) = \langle \mu , g \rangle ,
\]

where \(g(x) = \int f(sx) d\lambda_0(s)\) . Now, \(g\) is continuous (Ch. VII, §1, No.~1, Lemma 1). On the other hand, let \(S\) be the support of \(\lambda_0\) and \(K\) that of \(f\) . The conditions \(sx \in K\) and \(s \in S\) imply \(x \in S^{-1}K\) ; therefore the support of \(g\) is contained in \(S^{-1}K\) , so that \(g \in \mathcal{K}(X)\) . Then \(\langle \lambda_0 * \mu, f \rangle = \langle \mu, g \rangle\) is a vaguely continuous function of \(\mu\) , which proves our assertion.

\begin{enumerate}
\def\labelenumi{\arabic{enumi})}
\setcounter{enumi}{1}
\tightlist
\item
  Let \(\mu_0 \in \mathcal{M}(\mathrm{X})\) . The mapping \(\lambda \mapsto \lambda * \mu_0\) of \(\mathcal{C}'(\mathrm{G})\) into \(\mathcal{M}(\mathrm{X})\) is continuous for the topologies \(\sigma(\mathcal{C}'(\mathrm{G}), \mathcal{C}(\mathrm{G}))\) and \(\sigma(\mathcal{M}(\mathrm{X}), \mathcal{K}(\mathrm{X}))\) . For, let \(f \in \mathcal{K}(\mathrm{X})\) . Setting \(h(s) = \int f(sx) d\mu_0(x)\) , we have \(\langle f, \lambda * \mu_0 \rangle = \langle h, \lambda \rangle\) , and \(h \in \mathcal{C}(\mathrm{G})\) (Ch. VII, §1, No.~1, Lemma 1).
\end{enumerate}

PROPOSITION 13. --- The mapping \((s,\mu)\mapsto\boldsymbol{\gamma}(s)\mu\) of \(\mathbf{G}\times\mathcal{M}_{+}(\mathbf{X})\) into \(\mathcal{M}_{+}(\mathbf{X})\) is continuous when the set \(\mathcal{M}_{+}(\mathbf{X})\) of positive measures on X is equipped with the vague topology.

Since \(\pmb{\gamma}(s)\mu = \pmb{\gamma}(ss_0^{-1})\pmb{\gamma}(s_0)\mu\) , it follows from Remark 1 that it suffices to prove the continuity of the mapping under consideration at a point of the form \((e,\mu_0)\) with \(\mu_0\in \mathcal{M}_+(\mathrm{X})\) . Given a function \(f\in \mathcal{K}(\mathrm{X})\) and a number \(\varepsilon >0\) , it is thus a matter of showing that there exist a neighborhood U of \(e\) in G and a neighborhood W of \(\mu_0\) in \(\mathcal{M}_+(\mathrm{X})\) such that the relations \(s\in \mathrm{U}\) , \(\mu \in \mathrm{W}\) imply

\[
\left| \int f (s x) d \mu (x) - \int f (x) d \mu_ {0} (x) \right| \leqslant \varepsilon .\tag{6}
\]

Let \(\mathbf{V}\) be a compact neighborhood of the support \(\mathbf{K}\) of \(f\) in \(\mathbf{X}\) , and let \(\varphi \in \mathcal{K}_{+}(\mathbf{X})\) be such that \(\varphi(x) = 1\) on \(\mathbf{V}\) ; there exists a neighborhood \(\mathrm{W}_0\) of \(\mu_0\) in \(\mathcal{M}_{+}(\mathbf{X})\) such that \(a = \sup_{\mu \in \mathrm{W}_0} \mu(\mathrm{V})\) is finite: it suffices to take for \(\mathrm{W}_0\) the set of \(\mu \in \mathcal{M}_{+}(\mathbf{X})\) such that \(|\langle \varphi, \mu - \mu_0 \rangle| \leqslant 1\) . Since the mapping \((s, x) \mapsto sx\) is continuous, there is, on the other hand, a compact neighborhood \(\mathrm{U}_0\) of \(e\) in \(\mathbf{G}\) such that \(s\mathrm{K} \subset \mathrm{V}\) for all \(s \in \mathrm{U}_0\) ; the function \((s, x) \mapsto f(sx)\) is then uniformly continuous in \(\mathrm{U}_0 \times \mathrm{V}\) and so there is a neighborhood \(\mathrm{U} \subset \mathrm{U}_0\) of \(e\) such that \(|f(sx) - f(x)| \leqslant \varepsilon / 2a\) for all \(s \in \mathrm{U}\) and \(x \in \mathrm{V}\) . For \(s \in \mathrm{U}\) and \(\mu \in \mathrm{W}_0\) , we therefore have

\[
\left| \int f (s x) d \mu (x) - \int f (x) d \mu (x) \right| \leqslant \varepsilon / 2;
\]

if \(\mathbf{W} \subset \mathbf{W}_0\) is the neighborhood of \(\mu_0\) in \(\mathcal{M}_{+}(\mathbf{X})\) formed by the measures \(\mu \in \mathbf{W}_0\) such that \(|\int f(x) d\mu(x) - \int f(x) d\mu_0(x)| \leqslant \varepsilon/2\) , \(\mathbf{U}\) and \(\mathbf{W}\) meet the requirements.

\section{CONVOLUTION OF MEASURES AND FUNCTIONS}

\subsection{Convolution of a measure and a function}

Let X be a locally compact space on which a locally compact group G operates on the left continuously. Let \(\beta\) be a positive measure on X, quasi-invariant under G. Let \(\chi\) be a function \textgreater0 on \(G \times X\) , measurable for every measure on \(\mathbf{G} \times \mathbf{X}\) , and such that, for every \(s \in \mathbf{G}\) , \(\chi(s^{-1}, \cdot)\) is a density of \(\pmb{\gamma}(s)\beta\) with respect to \(\beta\) :

\[
\pmb {\gamma} (s) \beta = \chi (s ^ {- 1}, \cdot) \cdot \beta ,\tag{1}
\]

which, with the conventions of Ch. VII, §1, No.~1, may be written:

\[
d \beta (s x) = \chi (s, x) d \beta (x).\tag{1'}
\]

These data will remain fixed in Nos. 1, 2, 3 (an exception being made in Remark 2 of No.~2).

Recall (§2, No.~5) that if \(\chi\) is continuous and \(\beta\) has support \(X\) , then \(\chi\) is a multiplier.

Let \(f\) be a locally \(\beta\) -integrable complex function on \(X\) , and let \(\mu\) be a measure on \(G\) . For every \(s \in G\) , the measure \(\boldsymbol{\gamma}(s)(f \cdot \beta)\) has base \(\beta\) since \(\beta\) is quasi-invariant. Therefore, if \(\mu\) and \(f \cdot \beta\) are convolvable, then \(\mu * (f \cdot \beta)\) has base \(\beta\) (§3, No.~2, Prop. 10).

DEFINITION 1. --- If \(\mu\) and \(f \cdot \beta\) are convolvable, \(\mu\) and \(f\) are said to be convolvable relative to \(\beta\) . Every density of \(\mu * (f \cdot \beta)\) with respect to \(\beta\) is called a convolution product of \(\mu\) and \(f\) relative to \(\beta\) and is denoted \(\mu * ^{\beta} f\) .

One omits \(\beta\) when no confusion is possible. Convolution for several measures on G and a function on X is defined in an analogous manner.

The various convolution products of \(\mu\) and f are equal locally \(\beta\) -almost everywhere. If \(\beta\) has support X and if there exists a convolution product of \(\mu\) and f that is continuous, then the latter is uniquely determined; it is then called the convolution product of \(\mu\) and f relative to \(\beta\) .

Let \(s \in G\) and let \(f\) be a locally \(\beta\) -integrable complex function on \(X\) . Then \(\varepsilon_s\) and \(f\) are convolvable, and

\[
\varepsilon_ {s} * (f \cdot \beta) = \boldsymbol {\gamma} (s) (f \cdot \beta) = (\boldsymbol {\gamma} (s) f) \cdot (\boldsymbol {\gamma} (s) \beta) = (\boldsymbol {\gamma} (s) f) \cdot \chi (s ^ {- 1}, \cdot) \cdot \beta ,
\]

therefore

\[
\left(\varepsilon_ {s} * f\right) (x) = \chi \left(s ^ {- 1}, x\right) f \left(s ^ {- 1} x\right) = \left(\gamma_ {\chi} (s) f\right) (x)\tag{2}
\]

locally \(\beta\) -almost everywhere.

Lemma 1. --- Let \(\mu\) be a measure on \(G\) . Then \(\chi\) is locally \((\mu \otimes \beta)\) -integrable, and the image of \(\mu \otimes \beta\) under the homeomorphism \((s, x) \mapsto (s, s^{-1}x)\) of \(G \times X\) onto \(G \times X\) is \(\chi \cdot (\mu \otimes \beta)\) .

We may suppose that \(\mu \geqslant 0\) . Let \(\mathbf{F} \in \mathcal{K}_{+}(\mathbf{G} \times \mathbf{X})\) . Then

\[
\begin{array}{l} \iint \mathrm{F} (s, s ^ {- 1} x) d \mu (s) d \beta (x) = \int d \mu (s) \int \mathrm{F} (s, s ^ {- 1} x) d \beta (x) \\ = \int d \mu (s) \int \mathrm{F} (s, x) d (\boldsymbol {\gamma} (s ^ {- 1}) \beta) (x) = \int d \mu (s) \int \mathrm{F} (s, x) \chi (s, x) d \beta (x). \end{array}
\]

Now, the function \((s,x)\mapsto\mathrm{F}(s,x)\chi(s,x)\) has compact support and is \((\mu\otimes\beta)\) -measurable. By Ch. V, §8, No.~3, Prop. 7, the preceding equality proves that this function is \((\mu\otimes\beta)\) -integrable and that

\[
\iint \mathrm{F} (s, s ^ {- 1} x) d \mu (s) d \beta (x) = \iint \mathrm{F} (s, x) \chi (s, x) d \mu (s) d \beta (x).
\]

This proves at the same time both assertions of Lemma 1.

PROPOSITION 1. --- Let \(\mu\) be a measure on G, \(f\) a locally \(\beta\) -integrable complex function on X. Suppose that the function \(s \mapsto f(s^{-1}x)\chi(s^{-1}, x)\) is essentially \(\mu\) -integrable except for a locally \(\beta\) -negligible set of values of \(x\) , and that the function \(x \mapsto \int |f(s^{-1}x)| \chi(s^{-1}, x) d|\mu|(s)\) , defined locally almost everywhere for \(\beta\) , is locally \(\beta\) -integrable. Then \(\mu\) and \(f\) are convolvable.

We may assume that \(f \geqslant 0\) and \(\mu \geqslant 0\) . Let \(h \in \mathcal{H}_{+}(X)\) . We are to prove that the function \((s, x) \mapsto h(sx)\) is essentially integrable for \(\mu \otimes (f \cdot \beta) = (1 \otimes f) \cdot (\mu \otimes \beta)\) (Ch. V, §8, No.~5, Prop. 10), that is, that \(\iint^{\bullet} h(sx)f(x)d\mu(s)d\beta(x) < +\infty\) (Ch. V, §5, No.~3, Prop. 3); it will clearly suffice to prove that there exists an \(a > 0\) such that for every compact subset \(K\) of \(G\) ,

\[
\iint^ {\bullet} h (s x) f (x) \varphi_ {K} (s) d \mu (s) d \beta (x) \leqslant a.
\]

By Lemma 1,

\[
\begin{array}{r l} \iint^ {\bullet} h (s x) f (x) \varphi_ {K} (s) d \mu (s) d \beta (x) \\ & = \iint^ {*} h (x) f (s ^ {- 1} x) \varphi_ {K} (s) \chi (s ^ {- 1}, x) d \mu (s) d \beta (x). \end{array}
\]

Now, the function \((s,x)\mapsto h(x)f(s^{-1}x)\varphi_{\mathrm{K}}(s)\chi(s^{-1},x)\) is \((\mu\otimes\beta)\) -measurable (Lemma 1) and has compact support. The preceding expression is therefore equal (Ch. V, §8, No.~3, Prop. 7) to

\[
\begin{array}{r l} & {\int^ {*} h (x) d \beta (x) \int^ {*} f (s ^ {- 1} x) \varphi_ {\mathrm{K}} (s) \chi (s ^ {- 1}, x) d \mu (s)} \\ & {\qquad \leqslant (\sup h) \int_ {\mathrm{S}} ^ {*} d \beta (x) \int^ {\bullet} f (s ^ {- 1} x) \chi (s ^ {- 1}, x) d \mu (s),} \end{array}
\]

where S denotes the support of h. Whence the proposition.

PROPOSITION 2. --- Let \(\mu\) be a measure on G, \(f\) a locally \(\beta\) -integrable complex function on X. Assume that one of the following conditions is satisfied:

\begin{enumerate}
\def\labelenumi{(\roman{enumi})}
\item
  \(f\) and \(\chi\) are continuous;
\item
  G operates properly in X and \(f\) is zero on the complement of a countable union of compact sets;
\item
  \(\mu\) is carried by a countable union of compact sets.
\end{enumerate}

If \(\mu\) and \(f\) are convolvable, then the function \(s \mapsto f(s^{-1}x)\chi(s^{-1}, x)\) is essentially \(\mu\) -integrable except for a locally \(\beta\) -negligible set of values of \(x\) , and one has, locally almost \(\beta\) -everywhere,

\[
(\mu * ^ {\beta} f) (x) = \int_ {G} f (s ^ {- 1} x) \chi (s ^ {- 1}, x) d \mu (s) = \int_ {G} (\gamma_ {\chi} (s) f) (x) d \mu (s).\tag{3}
\]

Let \(h \in \mathcal{K}(\mathrm{X})\) . Since \(\mu\) and \(f\) are convolvable, the function \((s, x) \mapsto h(sx)f(x)\) is essentially \((\mu \otimes \beta)\) -integrable. By Lemma 1, the function \((s, x) \mapsto h(x)f(s^{-1}x)\chi(s^{-1}, x)\) is essentially \((\mu \otimes \beta)\) -integrable. Under hypothesis (i) or (ii) of the statement, one then deduces that this function is \((\mu \otimes \beta)\) -integrable; for, in the first case it is continuous and one applies Prop. 3 of Ch. V, §1, No.~1, and in the second case it is zero outside a countable union of compact sets, and one applies Prop. 7, 2) of No.~2, loc. cit. By the Lebesgue--Fubini theorem,

\[
\begin{array}{c} \iint h (s x) d \mu (s) d (f \cdot \beta) (x) = \iint h (x) f (s ^ {- 1} x) \chi (s ^ {- 1}, x) d \mu (s) d \beta (x) \\ = \int h (x) d \beta (x) \int f (s ^ {- 1} x) \chi (s ^ {- 1}, x) d \mu (s), \end{array}
\]

the function \(x \mapsto g(x) = \int f(s^{-1}x)\chi (s^{-1},x)d\mu (s)\) being moreover locally \(\beta\) -integrable. One thus sees that

\[
\langle h, \mu * (f \cdot \beta) \rangle = \langle h, g \cdot \beta \rangle ,
\]

whence \(g = \mu *^{\beta} f\) .

Suppose now that \(\mu\) is carried by the union \(S\) of a sequence of compact sets. The function

\[
(s, x) \mapsto h (x) f (s ^ {- 1} x) \chi (s ^ {- 1}, x) \varphi_ {\mathrm{S}} (s)
\]

is essentially \((\mu\otimes\beta)\) -integrable, and zero outside a countable union of compact sets, hence \((\mu\otimes\beta)\) -integrable. Since \(\mu=\varphi_{S}\cdot\mu\) , the argument concludes as before.

Remark. --- The hypothesis (iii) of Prop. 2 is satisfied notably when \(\mu\) is bounded. For, for every \(n > 0\) , there then exists a compact subset \(\mathbf{K}_n\) of \(\mathbf{G}\) such that

\[
| \mu | (\mathrm{G} - \mathrm{K} _ {n}) \leqslant \frac {1}{n}
\]

(Ch. IV, §4, No.~7), and \(\mu\) is carried by the union of the \(K_{n}\) . More generally, let \(\rho\) be a lower semi-continuous finite function \(>0\) on \(G\) such that \(\rho(st) \leqslant \rho(s)\rho(t)\) ; if \(\mu \in \mathcal{M}^{\rho}\) , the hypothesis (iii) is satisfied; for, \(\rho \cdot \mu\) is bounded, and \(\mu\) is carried by the same subsets as \(\rho \cdot \mu\) since, on every compact subset of \(G\) , \(\rho\) is bounded below by a constant \(>0\) .

\subsection{Examples of convolvable measures and functions}

In Props. 3 and 4, \(\mathcal{C}'(\mathrm{G})\) and \(\mathcal{M}(\mathrm{G})\) are equipped with the topology of compact convergence in \(\mathcal{C}(\mathrm{G})\) and \(\mathcal{K}(\mathrm{G})\) , respectively.

PROPOSITION 3. --- Assume \(\chi\) continuous. Let \(\mu \in \mathcal{C}'(\mathbf{G})\) , \(f \in \mathcal{C}(\mathbf{X})\) . Then:

\begin{enumerate}
\def\labelenumi{(\roman{enumi})}
\item
  \(\mu\) and \(f\) are convolvable relative to \(\beta\) .
\item
  Formula (3) of No.~1 defines for every \(x \in \mathbf{X}\) a convolution product \(\mu *^{\beta}f\) that is continuous and is none other than the element \(\pmb{\gamma}_{\chi}(\mu)f\) defined by the continuous representation \(\pmb{\gamma}_{\chi}\) of \(\mathbf{G}\) in \(\mathcal{C}(\mathbf{X})\) ; moreover, the mapping \((\mu, f) \mapsto \mu *^{\beta}f\) is hypocontinuous relative to the equicontinuous subsets of \(\mathcal{C}'(\mathbf{G})\) and the compact subsets of \(\mathcal{C}(\mathbf{X})\) .
\item
  If in addition \(f \in \mathcal{K}(\mathrm{X})\) , then the product \(\mu *^{\beta}f\) of (ii) belongs to \(\mathcal{K}(\mathrm{X})\) and the mapping \((\mu, f) \mapsto \mu *^{\beta}f\) is hypocontinuous relative to the equicontinuous subsets of \(\mathcal{C}'(\mathrm{G})\) and the compact subsets of \(\mathcal{K}(\mathrm{X})\) .
\end{enumerate}

We know that \(\mu\) and \(f\) are convolvable (§3, No.~2, Prop. 8 (i)). On the other hand, with the notations of §2, we have

\[
\boldsymbol {\gamma} _ {\chi} (\mu) f = \int \left(\boldsymbol {\gamma} _ {\chi} (s) f\right) d \mu (s) \in \mathcal {C} (\mathrm{X})
\]

since \(\mathcal{C}(\mathbf{X})\) is complete. In particular, for every \(x\in \mathbf{X}\) ,

\[
\left(\boldsymbol {\gamma} _ {\chi} (\mu) f\right) (x) = \int \left(\boldsymbol {\gamma} _ {\chi} (s) f\right) (x) d \mu (s).
\]

This, combined with Prop. 2 (i), and §2, No.~6, proves (ii). Finally, if \(f \in \mathcal{K}(\mathrm{X})\) then \(\mu * (f \cdot \beta)\) has compact support (§3, No.~2, Prop. 9), therefore \(\mu *^{\beta} f \in \mathcal{K}(\mathrm{X})\) . For, let us consider the continuous representation \(U\) of \(\mathbf{G}\) in the completion \(\mathcal{K}(\mathrm{X})^{\wedge}\) obtained by extending by continuity the continuous operators \(\pmb{\gamma}_{\chi}(s)\) in \(\mathcal{K}(\mathbf{X})\) (§2, No.~1, Remark 3). Let S be the support of \(\mu\) . The functions \(\pmb{\gamma}_{\chi}(s)f\) , for \(s \in \mathbf{S}\) , have their support contained in a fixed compact set K. The set \(\mathcal{K}(\mathbf{X}, \mathbf{K})\) is a complete linear subspace of \(\mathcal{K}(\mathbf{X})\) . Therefore \(U(\mu)f \in \mathcal{K}(\mathbf{X})\) . One sees as before that \(U(\mu)f = \mu^{*\beta}f\) , and (iii) again follows from §2, No.~6.

PROPOSITION 4. --- Assume that G operates properly in X and that \(\chi\) is continuous. Let \(\mu \in \mathcal{M}(\mathrm{G})\) and \(f \in \mathcal{K}(\mathrm{X})\) .

\begin{enumerate}
\def\labelenumi{(\roman{enumi})}
\item
  \(\mu\) and \(f\) are convolvable relative to \(\beta\) .
\item
  Formula (3) of No.~1 defines for every \(x \in \mathbf{X}\) a convolution product \(\mu *^{\beta}f\) that is continuous.
\item
  The mapping \((\mu, f) \mapsto \mu *^{\beta}f\) of \(\mathcal{M}(\mathrm{G}) \times \mathcal{K}(\mathrm{X})\) into \(\mathcal{C}(\mathrm{X})\) is hypocontinuous relative to the bounded subsets of \(\mathcal{M}(\mathrm{G})\) and the compact subsets of \(\mathcal{K}(\mathrm{X})\) that are contained in some subspace \(\mathcal{K}(\mathrm{X}, \mathrm{L})\) (where \(\mathrm{L}\) is a variable compact subset of \(\mathrm{X}\) ).
\end{enumerate}

We know that \(\mu\) and \(f\) are convolvable (§3, No.~2, Prop. 8 (ii)), and it is clear that the integrals occurring in (3) exist for every \(x \in \mathbf{X}\) . Let \(\mathbf{K}\) and \(\mathbf{L}\) be two compact subsets of \(\mathbf{X}\) . There exists a compact subset \(\mathbf{H}\) of \(\mathbf{G}\) such that the relations \(x \in \mathbf{K}\) and \(s^{-1}x \in \mathbf{L}\) imply \(s \in \mathbf{H}\) ; let \(\varphi \in \mathcal{H}_{+}(\mathbf{G})\) with \(\varphi(s) = 1\) for \(s \in \mathbf{H}\) . Then, for \(f \in \mathcal{H}(\mathbf{X}, \mathbf{L})\) and \(x \in \mathbf{K}\) ,

\[
\begin{array}{c} \int f (s ^ {- 1} x) \chi (s ^ {- 1}, x) d \mu (s) = \int f (s ^ {- 1} x) \chi (s ^ {- 1}, x) \varphi (s) d \mu (s) \\ = ((\varphi \cdot \mu) * ^ {\beta} f) (x). \end{array}
\]

Consequently \(\int f(s^{-1}x)\chi(s^{-1},x)d\mu(s)\) is a continuous function of \(x\) and defines a convolution product \(\mu * ^ {\beta} f\in \mathcal{C}(\mathrm{X})\) . Moreover, the mapping \(\mu \mapsto \varphi \cdot \mu\) of \(\mathcal{M}(\mathrm{G})\) into \(\mathcal{C}'(\mathrm{G})\) is continuous for the topologies of compact convergence. Prop. 3 (iii) therefore implies that the mapping \((\mu ,f)\mapsto \mu * ^ {\beta} f\) of \(\mathcal{M}(\mathrm{G})\times \mathcal{K}(\mathrm{X},\mathrm{L})\) into \(\mathcal{C}(\mathrm{X})\) is, for every compact subset L of X, hypocontinuous relative to the compact subsets of \(\mathcal{K}(\mathrm{X},\mathrm{L})\) . In particular, the mapping \((\mu ,f)\mapsto \mu * ^ {\beta} f\) of \(\mathcal{M}(\mathrm{G})\times \mathcal{K}(\mathrm{X})\) into \(\mathcal{C}(\mathrm{X})\) is separately continuous. Since \(\mathcal{K}(\mathrm{X})\) is barreled, this mapping is hypocontinuous relative to the bounded subsets of \(\mathcal{M}(\mathrm{G})\) (TVS, III, §5, No.~3, Prop. 6).

Remark 1. --- Under the hypotheses of Prop. 4, the mapping \(\mu \mapsto \mu *^{\beta}f\) of \(\mathcal{M}_{+}(\mathrm{G})\) into \(\mathcal{C}(\mathrm{X})\) is continuous when \(\mathcal{M}_{+}(\mathrm{G})\) is equipped with the vague topology, for every \(f\in \mathcal{K}(\mathrm{X})\) . For, let K be a compact subset of X, S the (compact) support of \(f\) ; since G operates properly in X, the set of \(s\in \mathrm{G}\) for which there exists an \(x\in \mathrm{K}\) with \(s^{-1}x\in \mathrm{S}\) is a compact subset L of G (GT, III, §4, No.~5, Th. 1). Let \(\varepsilon\) be a number \(>0\) , \(\varphi\) a function in \(\mathcal{K}_{+}(\mathrm{G})\) equal to 1 on the compact set L, \(\mu_0\) an element of \(\mathcal{M}_{+}(\mathrm{G})\) ; the set \(\mathrm{W}_0\) of measures \(\mu \in \mathcal{M}_{+}(\mathrm{G})\) such that

\[
\left| \int \varphi (s) d \mu (s) - \int \varphi (s) d \mu_ {0} (s) \right| \leqslant \varepsilon
\]

is a neighborhood of \(\mu_0\) in \(\mathcal{M}_{+}(\mathrm{G})\) . On the other hand, the function \((s,x)\mapsto f(s^{-1}x)\chi (s^{-1},x)\) is uniformly continuous on \(\mathbf{L}\times \mathbf{K}\) , therefore there exists a finite number of points \(x_{i}\in \mathbf{K}\) ( \(1\leqslant i\leqslant n\) ) such that for every \(x\in \mathbf{K}\) , there is an \(i\) for which

\[
\left| f (s ^ {- 1} x) \chi (s ^ {- 1}, x) - f (s ^ {- 1} x _ {i}) \chi (s ^ {- 1}, x _ {i}) \right| \leqslant \varepsilon
\]

for all \(s \in \mathbf{L}\) . Since \(\mu(\mathbf{L}) \leqslant \int \varphi(s) d\mu_0(s) + \varepsilon\) for all \(\mu \in \mathbf{W}_0\) , also

\[
\begin{array}{r l} \Big | \int f (s ^ {- 1} x) \chi (s ^ {- 1}, x) d \mu (s) - \int f (s ^ {- 1} x _ {i}) \chi (s ^ {- 1}, x _ {i}) d \mu (s) \Big | & \\ & \leqslant \varepsilon \Big (\int \varphi (s) d \mu_ {0} (s) + \varepsilon \Big) \end{array}
\]

for every x satisfying the preceding inequality and every \(\mu \in W_{0}\) . Now let W be the neighborhood of \(\mu_{0}\) in \(\mathcal{M}_{+}(G)\) formed by the measures \(\mu \in W_{0}\) such that

\[
\left| \int f (s ^ {- 1} x _ {i}) \chi (s ^ {- 1}, x _ {i}) d \mu (s) - \int f (s ^ {- 1} x _ {i}) \chi (s ^ {- 1}, x _ {i}) d \mu_ {0} (s) \right| \leqslant \varepsilon
\]

for \(1 \leqslant i \leqslant n\) . It is clear that for every measure \(\mu \in \mathbf{W}\) and every \(x \in \mathbf{K}\) ,

\[
\begin{array}{r l} & {\Big | \int f (s ^ {- 1} x) \chi (s ^ {- 1}, x) d \mu (s) - \int f (s ^ {- 1} x) \chi (s ^ {- 1}, x) d \mu_ {0} (s) \Big |} \\ & {\qquad \leqslant \varepsilon \Big (2 \int \varphi (s) d \mu_ {0} (s) + 2 \varepsilon + 1 \Big),} \end{array}
\]

and since \(\varepsilon\) is arbitrary, this proves our assertion.

PROPOSITION 5. --- Assume \(\chi\) a continuous multiplier and each function \(\chi(s,\cdot)\) bounded.

\begin{enumerate}
\def\labelenumi{(\roman{enumi})}
\item
  The function \(s \mapsto \rho(s) = \sup_{x \in X} \chi(s^{-1}, x)\) on \(G\) is lower semi-continuous \(> 0\) and satisfies \(\rho(st) \leqslant \rho(s)\rho(t)\) for all \(s, t\) in \(G\) .
\item
  Let \(\mu \in \mathcal{M}^{\rho}(\mathrm{G})\) and \(f \in \mathrm{L}^{\infty}(\mathrm{X},\beta)\) . Then \(\mu\) and \(f\) are convolvable and \(\mu *^{\beta} f\) is given locally almost everywhere by the formula (3) of No.~1. One has \(\mu *^{\beta} f \in \mathrm{L}^{\infty}(\mathrm{X},\beta)\) , and \(\| \mu *^{\beta} f \|_{\infty} \leqslant \| \mu \|_{\rho} \| f \|_{\infty}\) .
\item
  If, moreover, \(f \in \mathcal{C}^{\infty}(\mathrm{X})\) (resp. \(\overline{\mathcal{K}(\mathrm{X})}\) ), then formula (3) of No.~1 defines for every \(x\) a convolution product \(\mu *^{\beta}f\) that belongs to \(\mathcal{C}^{\infty}(\mathrm{X})\) (resp. \(\overline{\mathcal{K}(\mathrm{X})}\) ).
\item
  If \(f \in \overline{\mathcal{K}(\mathrm{X})}\) , then the convolution product \(\mu *^{\beta}f\) defined by (3) is none other than the element \(\gamma_{\chi}(\mu)f\) defined by the continuous representation \(\gamma_{\chi}\) of \(G\) in \(\overline{\mathcal{K}(\mathrm{X})}\) .
\end{enumerate}

The identity \(\chi(st,x) = \chi(s,tx)\chi(t,x)\) implies at once that \(\rho(st) \leqslant \rho(s)\rho(t)\) . On the other hand, \(\rho\) is lower semi-continuous, being the upper envelope of continuous functions.

Let \(\mu \in \mathcal{M}^{\rho}(\mathrm{G})\) . By Prop. 1 of No.~1, \(\mu\) and 1 are convolvable; Prop. 2 (i) shows that \((|\mu|^{*\beta}1)(x) \leqslant \int_{\mathrm{G}} \rho(s)d|\mu|(s)\) locally \(\beta\) -almost everywhere. Therefore, if \(f\) is \(\beta\) -measurable and \(|f| \leqslant 1\) , then \(\mu\) and \(f\) are convolvable and \(\mathrm{N}_{\infty}(\mu^{*\beta}f) \leqslant \int \rho(s)d|\mu|(s)\) . Moreover, \(\mu^{*\beta}f\) is given locally almost everywhere by formula (3) of No.~1, because condition (iii) of Prop. 2 of No.~1 is satisfied. This implies (ii).

Suppose \(f\) continuous and bounded by 1 in absolute value. It is clear that the integrals occurring in (3) exist for all \(x \in X\) . Let us show that they depend continuously on \(x\) . We can suppose \(\mu \geqslant 0\) . Let \(x_0 \in X\) and \(\varepsilon > 0\) . Let \(K\) be a compact subset of \(G\) such that \(\int_{G - K} \rho(s) d\mu(s) \leqslant \varepsilon\) . There exists a neighborhood \(V\) of \(x_0\) in \(X\) such that \(x \in V\) implies

\[
\left| f \left(s ^ {- 1} x\right) \chi \left(s ^ {- 1}, x\right) - f \left(s ^ {- 1} x _ {0}\right) \chi \left(s ^ {- 1}, x _ {0}\right) \right| \leqslant \varepsilon / \mu (\mathrm{K})
\]

for \(s \in \mathbf{K}\) . Then, for \(x \in \mathbf{V}\) ,

\[
\begin{array}{r l} & {\Big | \int f (s ^ {- 1} x) \chi (s ^ {- 1}, x) d \mu (s) - \int f (s ^ {- 1} x _ {0}) \chi (s ^ {- 1}, x _ {0}) d \mu (s) \Big |} \\ & {\qquad \leqslant 2 \int_ {\mathrm{G-K}} \rho (s) d \mu (s) + \int_ {\mathrm{K}} \frac {\varepsilon}{\mu (\mathrm{K})} d \mu (s) \leqslant 3 \varepsilon ,} \end{array}
\]

whence our assertion. Suppose that in addition \(f \in \overline{\mathcal{H}(\mathrm{X})}\) . Let \(\mathbf{H}\) be a compact subset of \(\mathbf{X}\) such that \(|f(y)| \leqslant \varepsilon\) for \(y \notin \mathbf{H}\) . Let \(x \notin \mathbf{KH}\) . Then \(s^{-1}x \notin \mathbf{H}\) for \(s \in \mathbf{K}\) , therefore

\[
\begin{array}{r l} \Big | \int_ {\mathrm{G}} f (s ^ {- 1} x) \chi (s ^ {- 1}, x) d \mu (s) \Big | & \leqslant \int_ {\mathrm{G} - \mathrm{K}} \rho (s) d \mu (s) + \int_ {\mathrm{K}} \varepsilon \rho (s) d \mu (s) \\ & \leqslant \varepsilon \Big (1 + \int_ {\mathrm{G}} \rho (s) d \mu (s) \Big), \end{array}
\]

which completes the proof of (iii).

Finally, if \(f \in \overline{\mathcal{K}(\mathbf{X})}\) then, since \(\varepsilon_x \in \mathcal{M}^1(\mathbf{X})\) for all \(x \in \mathbf{X}\) , we have

\[
\left(\boldsymbol {\gamma} _ {\chi} (\mu) f\right) (x) = \int \left(\boldsymbol {\gamma} _ {\chi} (s) f\right) (x) d \mu (s),
\]

thus \(\gamma_{\chi}(\mu)f\) is the convolution product \(\mu *^{\beta}f\) defined by (3).

PROPOSITION 6. --- Assume \(\chi\) a continuous multiplier and each function \(\chi(s, \cdot)\) bounded. Let \(\rho(s) = \sup_{x \in X} \chi(s^{-1}, x)\) . Let \(p\) and \(q\) be two conjugate exponents ( \(1 \leqslant p < +\infty\) ). Let \(\mu \in \mathcal{M}^{\rho^{1/q}}(G)\) and \(f \in L^{p}(X, \beta)\) . Then:

\begin{enumerate}
\def\labelenumi{(\roman{enumi})}
\item
  \(\mu\) and \(f\) are convolvable;
\item
  the convolution product \(\mu * ^ {\beta} f\) is given locally \(\beta\) -almost everywhere by the formula (3), and is equal locally \(\beta\) -almost everywhere to a function \(g\in \mathbf{L}^p (\mathbf{X},\beta)\) such that \(\| g\| _p\leqslant \| \mu \|_{\rho^{1 / q}}\| f\| _p;\) (iii) \(g\) is equal to the element \(\gamma_{\chi}(\mu)f\) defined by the continuous representation \(\gamma_{\chi}\) of \(\mathbf{G}\) in \(\mathbf{L}^p (\mathbf{X},\beta)\) .
\end{enumerate}

We have

\[
\int^ {*} \| \gamma_ {\chi} (s) f \| _ {p} d | \mu | (s) \leqslant \left(\int^ {*} \rho (s) ^ {1 / q} d | \mu | (s)\right) \| f \| _ {p} <   + \infty
\]

by §2, No.~5, formula (5). On the other hand, the mapping \(s \mapsto \pmb{\gamma}_{\chi}(s)f\) of G into \(\mathrm{L}^p (\mathbf{X},\beta)\) is continuous (§2, No.~5, Prop. 9). Therefore this mapping is \(\mu\) -integrable. Let

\[
g = \int_ {\mathrm{G}} \left(\boldsymbol {\gamma} _ {\chi} (s) f\right) d \mu (s) \in \mathrm{L} ^ {p} (\mathrm{X}, \beta).
\]

We have \(\| g\| _p\leqslant \left(\int \rho^{1 / q}(s)d|\mu |(s)\right)\| f\| _p\) . Applying the preceding remarks to \(|f|\) , one sees that the mapping \(s\mapsto \varepsilon_s*\mid f|\) of G into \(\mathbf{L}^p (\mathbf{X},\beta)\) is \(\mu\) -integrable, therefore that, for every \(h\in \mathcal{K}(\mathbf{X})\) , the mapping \(s\mapsto \langle h,\varepsilon_s*(|f|\cdot \beta)\rangle\) is \(\mu\) -integrable. Prop. 7 of §1, No.~5 then proves that \(\mu\) and \(f\cdot \beta\) are convolvable. Moreover,

\[
\begin{array}{r l} \int_ {\mathrm{X}} g (x) h (x) d \beta (x) & = \int_ {\mathrm{G}} d \mu (s) \int_ {\mathrm{X}} (\boldsymbol {\gamma} _ {\chi} (s) f) (x) h (x) d \beta (x) \\ & = \int_ {\mathrm{G}} \langle h, \varepsilon_ {s} * (f \cdot \beta) \rangle d \mu (s), \end{array}
\]

and this last integral is equal to \(\langle h, \mu * (f \cdot \beta) \rangle\) by Prop. 7 of §1, No.~5. One therefore sees that g is a convolution product of \(\mu\) and f.~This convolution product is given locally \(\beta\) -almost everywhere by (3), by Prop. 2 and the Remark that follows it.

By an abuse of notation, it is often one of the functions g of the statement that is denoted \(\mu * ^ {\beta} f\) , which permits writing

\[
\left\| \mu * ^ {\beta} f \right\| _ {p} \leqslant \left\| \mu \right\| _ {\rho^ {1 / q}} \| f \| _ {p}.
\]

If X is countable at infinity, this style of notation is, moreover, entirely justified.

COROLLARY. --- Under the hypotheses of Prop. 6, the mapping \((\mu, f) \mapsto \mu *^{\beta} f\) defines on \(\mathrm{L}^{p}(\mathrm{X}, \beta)\) the structure of a left module over \(\mathcal{M}^{\rho^{1/q}}(\mathrm{G})\) ( \(1 \leqslant p \leqslant +\infty\) ).

This follows from Props. 5 and 6 and the associativity of the convolution product.

Remark 2. --- Let X be a locally compact space on which a locally compact group G operates continuously on the right by \((x,s)\mapsto xs\) . Let \(\beta\) be a positive measure on X. Let \(\chi\) be a function \textgreater0 on \(G\times X\) , measurable for every measure on \(G\times X\) , such that \(\boldsymbol{\delta}(s)\boldsymbol{\beta}=\chi(s,\cdot)\cdot\boldsymbol{\beta}\) for every \(s\in G\) . Let f be a locally \(\beta\) -integrable function on X and let \(\mu\) be a measure on G. If \(f\cdot\beta\) and \(\mu\) are convolvable (for the mapping \((x,s)\mapsto xs\) of \(X\times G\) into X), then \((f\cdot\beta)*\mu\) has base \(\beta\) . We then say that f and \(\mu\) are convolvable relative to \(\beta\) ; every density of \((f\cdot\beta)*\mu\) with respect to \(\beta\) is called a convolution product of f and \(\mu\) relative to \(\beta\) , and is denoted \(f * ^ {\beta} \mu\) or simply \(f*\mu\) .

Let \(G^{0}\) be the group opposite to G. By \((s,x)\mapsto xs\) , \(G^{0}\) operates continuously on the left in X. To say that f and \(\mu\) are convolvable in the foregoing sense is equivalent to saying that \(\mu\) and f are convolvable for \(G^{0}\) operating on the left in X; and the convolution products \(f * ^ {\beta} \mu\) are none other than the convolution products \(\mu * ^ {\beta} f\) for \(G^{0}\) operating on the left in X. On the other hand, for \(s\in G^{0}\) one has \(\pmb{\gamma}(s)\beta = \chi(s^{-1},\cdot)\cdot \beta\) . The results of Nos. 1 and 2 may then be translated immediately into results concerning the products \(f * ^ {\beta} \mu\) . In particular:

\begin{enumerate}
\def\labelenumi{\arabic{enumi})}
\tightlist
\item
  If \(s \in \mathbf{G}\) and \(f\) is locally \(\beta\) -integrable, then \(f\) and \(\varepsilon_s\) are convolvable and
\end{enumerate}

\[
(f * \varepsilon_ {s}) (x) = \chi (s ^ {- 1}, x) f (x s ^ {- 1}).\tag{4}
\]

\begin{enumerate}
\def\labelenumi{\arabic{enumi})}
\setcounter{enumi}{1}
\tightlist
\item
  If \(f\) and \(\mu\) are convolvable and if one of the conditions (i), (ii), (iii) of Prop. 2 is fulfilled, then \(f *^{\beta} \mu\) is given locally \(\beta\) -almost everywhere by
\end{enumerate}

\[
(f * ^ {\beta} \mu) (x) = \int_ {G} f (x s ^ {- 1}) \chi (s ^ {- 1}, x) d \mu (s).\tag{5}
\]

We leave to the reader the task of translating the other statements. Note that if \(\chi\) is continuous and \(\beta\) has support X, then

\[
\chi (t s, x) = \chi (s, x t) \chi (t, x) \quad (x \in \mathrm{X}; s, t \text {in} \mathrm{G}).\tag{6}
\]

\subsection{Convolution and transposition}

Let us return to the hypotheses and notations of the beginning of No.~1, but let us assume in addition that \(\beta\) is relatively invariant with multiplier \(\chi\) ; \(\chi\) is therefore a continuous function on G.

PROPOSITION 7. --- Let \(f\) be a locally \(\beta\) -integrable function on \(X\) , \(\nu\) a measure on \(X\) , and \(\mu\) a measure on \(G\) . Assume that:

\begin{enumerate}
\def\labelenumi{(\roman{enumi})}
\item
  \(\mu\) and \(f\) are convolvable and formula (3) of No.~1 defines locally \(\beta\) -almost everywhere a convolution product \(\mu *^{\beta} f\) .
\item
  \(\chi \cdot \check{\mu}\) and \(\nu\) are convolvable.
\item
  The function \(g(s,x) = f(s^{-1}x)\chi (s^{-1})\) is \((\mu \otimes \nu)\) -integrable.
\end{enumerate}

Then \(f\) is essentially integrable for \((\chi \cdot \breve{\mu}) * \nu\) , the function \(\mu *^{\beta} f\) defined by (3) is \(\nu\) -integrable, and

\[
\nu (\mu * ^ {\beta} f) = \big ((\chi \cdot \check {\mu}) * \nu \big) (f).\tag{7}
\]

Since \(g(s, x)\) is integrable for \(\mu \otimes \nu\) , the function \(f(sx)\) is essentially integrable for \((\chi \cdot \breve{\mu}) \otimes \nu\) and \(f\) is essentially integrable for \((\chi \cdot \breve{\mu}) * \nu\) . By the Lebesgue-Fubini theorem, \(\mu *^{\beta} f = \int g(s, x) d\mu(s)\) is \(\nu\) -integrable and

\[
\begin{array}{l} \nu (\mu * ^ {\beta} f) = \iint f (s ^ {- 1} x) \chi (s ^ {- 1}) d \mu (s) d \nu (x) \\ = \iint f (s x) \chi (s) d \stackrel {\vee} {\mu} (s) d \nu (x) = \big ((\chi \cdot \stackrel {\vee} {\mu}) * \nu \big) (f). \end{array}
\]

Examples. --- 1) One can take \(f \in \mathcal{C}(\mathrm{X})\) , \(\nu \in \mathcal{C}'(\mathrm{X})\) and \(\mu \in \mathcal{C}'(\mathrm{G})\) by Prop. 3, and the Cor. of Prop. 5 of §1, No.~4. The formula (7) then means that the endomorphism \(\nu \mapsto (\chi \cdot \breve{\mu}) * \nu\) of \(\mathcal{C}'(\mathrm{X})\) is the transpose of the endomorphism \(f \mapsto \mu * f\) of \(\mathcal{C}(\mathrm{X})\) .

\begin{enumerate}
\def\labelenumi{\arabic{enumi})}
\setcounter{enumi}{1}
\item
  One can take \(f \in \mathcal{K}(\mathrm{X})\) , \(\nu \in \mathcal{M}(\mathrm{X})\) and \(\mu \in \mathcal{C}'(\mathrm{G})\) by Prop. 3, Prop. 8 of §3, No.~2, and the remark that the support of the continuous function \(g(s,x)\) intersects the support of \(\mu \otimes \nu\) in a compact set. The formula (7) then means that the endomorphism \(\nu \mapsto (\chi \cdot \breve{\mu}) * \nu\) of \(\mathcal{M}(\mathrm{X})\) is the transpose of the endomorphism \(f \mapsto \mu * f\) of \(\mathcal{K}(\mathrm{X})\) .
\item
  If G operates properly on X, one can take \(f \in \mathcal{K}(X)\) , \(\nu \in \mathcal{C}'(X)\) and \(\mu \in \mathcal{M}(G)\) by Prop. 4, Prop. 8 of §3, No.~2, and the same remark as in Example 2.
\end{enumerate}

PROPOSITION 8. --- Let \(f\) and \(g\) be two locally \(\beta\) -integrable functions on \(X\) and let \(\mu \in \mathcal{M}(G)\) . Assume that:

\begin{enumerate}
\def\labelenumi{(\roman{enumi})}
\item
  \(\mu\) and \(f\) are convolvable and the formula (3) of No.~1 defines locally \(\beta\) -almost everywhere a convolution product \(\mu *^{\beta} f\) .
\item
  \(\chi \cdot \breve{\mu}\) and \(g\) are convolvable and the formula (3) of No.~1 (with \(\mu\) replaced by \(\chi \cdot \breve{\mu}\) and \(f\) by \(g\) ) defines locally \(\beta\) -almost everywhere a convolution product \((\chi \cdot \breve{\mu}) *^{\beta} g\) .
\item
  There exists a function \(\psi\) on \(\mathbf{G}\) , equal locally \(\mu\) -almost everywhere to 1, such that the function
\end{enumerate}

\[
h (s, x) = g (x) f \left(s ^ {- 1} x\right) \chi \left(s ^ {- 1}\right) \psi (s)
\]

is \((\mu \otimes \beta)\) -integrable.

Then the functions \(g(x)\big((\mu * ^ {\beta} f)(x)\big)\) and \(f(x)\big(\big((\chi \cdot \breve{\mu}) * ^ {\beta} g\big)(x)\big)\) are essentially \(\beta\) -integrable, and

\[
\int f (x) \big (\big ((\chi \cdot \breve {\mu}) * ^ {\beta} g \big) (x) \big) d \beta (x) = \int g (x) \big ((\mu * ^ {\beta} f) (x) \big) d \beta (x).\tag{8}
\]

For, by (iii) and the Lebesgue--Fubini theorem, the function

\[
x \mapsto g (x) \int f (s ^ {- 1} x) \chi (s ^ {- 1}) \psi (s) d \mu (s)
\]

is \(\beta\) -integrable, and

\[
\begin{array}{l} \mathrm{I} = \iint f (s ^ {- 1} x) g (x) \chi (s ^ {- 1}) \psi (s) d \mu (s) d \beta (x) \\ = \int g (x) d \beta (x) \int f (s ^ {- 1} x) \chi (s ^ {- 1}) \psi (s) d \mu (s). \end{array}
\]

But \(\psi \cdot \mu = \mu\) , consequently

\[
\int f (s ^ {- 1} x) \chi (s ^ {- 1}) \psi (s) d \mu (s) = (\mu * ^ {\beta} f) (x)
\]

locally \(\beta\) -almost everywhere. This shows that the function

\[
x \mapsto g (x) \left(\left(\mu * ^ {\beta} f\right) (x)\right)
\]

is essentially \(\beta\) -integrable and that

\[
\mathrm{I} = \int g (x) \left(\left(\mu * ^ {\beta} f\right) (x)\right) d \beta (x).
\]

On the other hand, Lemma 1 shows that the function

\[
(s, x) \mapsto g (s x) f (x) \chi (s ^ {- 1}) \psi (s)
\]

is integrable for \((\chi \cdot \mu) \otimes \beta\) . Therefore the function \((s, x) \mapsto g(s^{-1}x)f(x)\psi(s^{-1})\) is integrable for \(\breve{\mu} \otimes \beta\) , and

\[
\begin{array}{l} \mathrm{I} = \iint g (s ^ {- 1} x) f (x) \psi (s ^ {- 1}) d \mu^ {\vee} (s) d \beta (x) \\ = \int f (x) d \beta (x) \int g (s ^ {- 1} x) \psi (s ^ {- 1}) d \mu^ {\vee} (s). \end{array}
\]

But \(\check{\psi} \cdot \check{\mu} = \check{\mu}\) and so \(\int g(s^{-1}x)\psi(s^{-1})d\check{\mu}(s) = ((\chi \cdot \check{\mu}) *^{\beta}g)(x)\) locally \(\beta\) -almost everywhere. This shows that the function

\[
x \mapsto f (x) \left(\left((\chi \cdot \stackrel {\vee} {\mu}) * ^ {\beta} g\right) (x)\right)
\]

is essentially \(\beta\) -integrable and that

\[
\mathrm{I} = \int f (x) \left(\left(\left(\chi \cdot \stackrel {\vee} {\mu}\right) * ^ {\beta} g\right) (x)\right) d \beta (x).
\]

This proves the proposition.

Examples. --- 4) One can take \(f \in \mathcal{C}(\mathrm{X})\) , \(g \in \mathcal{K}(\mathrm{X})\) and \(\mu \in \mathcal{C}'(\mathrm{G})\) (with \(\psi = 1\) ).

\begin{enumerate}
\def\labelenumi{\arabic{enumi})}
\setcounter{enumi}{4}
\item
  If \(\mathbf{G}\) operates properly on \(\mathbf{X}\) , one can take \(f \in \mathcal{K}(\mathbf{X})\) , \(g \in \mathcal{K}(\mathbf{X})\) and \(\mu \in \mathcal{M}(\mathbf{G})\) (with \(\psi = 1\) ).
\item
  One can take \(f \in \mathrm{L}^p(\mathrm{X},\beta)\) , \(g \in \mathrm{L}^q(\mathrm{X},\beta)\) and \(\mu \in \mathcal{M}^\rho(\mathrm{G})\) , where \(1 \leqslant p < +\infty\) , \(\frac{1}{p} + \frac{1}{q} = 1\) , \(\rho = \chi^{-1/q}\) . The conditions (i) and (ii) are satisfied by Props. 5 and 6. Let us prove (iii). We have seen that \(\mu\) is carried by a set \(S\) that is a countable union of compact sets. Let us take for \(\psi\) the characteristic function of \(S\) . The function \(h\) is \((\mu \otimes \beta)\) -measurable: for, the function \((s,x) \mapsto g(x)\chi(s^{-1})\psi(s)\) is so, as is the function \((s,x) \mapsto f(s^{-1}x)\) by Lemma 1. Moreover, \(g\) being zero outside a countable union of \(\beta\) -integrable sets, \(h\) is zero outside a countable union of \((\mu \otimes \beta)\) -integrable sets. We then have (Ch. V, §8, No.~3, Prop. 7):
\end{enumerate}

\[
\begin{array}{l} \mathrm{J} = \iint^ {*} | g (x) f (s ^ {- 1} x) | \chi (s ^ {- 1}) \psi (s) d | \mu | (s) d \beta (x) \\ = \int^ {*} | g (x) | d \beta (x) \int^ {*} | f (s ^ {- 1} x) | \chi (s ^ {- 1}) \psi (s) d | \mu | (s). \end{array}\tag{9}
\]

But since \(g\) (resp. \(\psi\) ) is zero outside a countable union of integrable sets, the upper integrals of the second member of (9) are equal to the essential upper integrals (Ch. V, §1, No.~2, Prop. 7). Now (Ch. V, §5, No.~3, Prop. 3)

\[
\int^ {\bullet} | f (s ^ {- 1} x) | \chi (s ^ {- 1}) \psi (s) d | \mu | (s) = \int^ {\bullet} | f (s ^ {- 1} x) | \chi (s ^ {- 1}) d | \mu | (s)
\]

since \(\mu = \psi \cdot \mu\) . By Prop. 6, this last integral is finite and is equal to \((|\mu| *^{\beta} |f|)(x)\) locally \(\beta\) -almost everywhere. Therefore

\[
J = \int^ {\bullet} | g (x) | \left(| \mu | * ^ {\beta} | f |\right) (x) d \beta (x),
\]

and J is finite since \(g \in \mathbf{L}^q\) and \(|\mu| *^\beta |f| \in \mathbf{L}^p\) (Prop. 6). Therefore \(h\) is \((\mu \otimes \beta)\) -integrable.

The formula (8) then means that the endomorphism \(g \mapsto (\chi \cdot \breve{\mu}) * g\) of \(\mathrm{L}^q(\mathrm{X},\beta)\) is, for \(\mu \in \mathcal{M}^\rho(\mathrm{G})\) , the transpose of the endomorphism \(f \mapsto \mu * f\) of \(\mathrm{L}^p(\mathrm{X},\beta)\) .

\subsection{Convolution of a measure and a function on a group}

Let G be a locally compact group. Throughout Nos. 4 and 5, we fix a relatively invariant positive measure \(\beta \neq 0\) on G; let \(\chi\) and \(\chi'\) be its left and right multipliers (recall that \(\chi' = \chi \Delta_{\mathrm{G}}\) ). If \(\mu\) is a measure on G and \(f\) is a locally \(\beta\) -integrable function on G, the convolvability of \(\mu\) and \(f\) and the products \(\mu * f\) (resp. the convolvability of \(f\) and \(\mu\) and the products \(f * \mu\) ) may be defined by regarding G as operating on itself on the left (resp. on the right) by translations. Let us make explicit in this situation some of the preceding results:

\begin{enumerate}
\def\labelenumi{\arabic{enumi})}
\tightlist
\item
  Let \(\mu\) be a measure on \(\mathbf{G}\) , \(f\) a locally \(\beta\) -integrable complex function on \(\mathbf{G}\) . Assume verified one of the following conditions:
\end{enumerate}

\begin{enumerate}
\def\labelenumi{(\roman{enumi})}
\item
  \(f\) is continuous;
\item
  \(f\) is zero on the complement of a countable union of compact sets;
\item
  \(\mu\) is carried by a countable union of compact sets.
\end{enumerate}

If \(\mu\) and \(f\) are convolvable then, locally \(\beta\) -almost everywhere,

\[
(\mu * f) (x) = \int_ {G} f (s ^ {- 1} x) \chi (s ^ {- 1}) d \mu (s).\tag{10}
\]

If \(f\) and \(\mu\) are convolvable then, locally \(\beta\) -almost everywhere,

\[
(f * \mu) (x) = \int_ {\mathrm{G}} f (x s ^ {- 1}) \chi^ {\prime} (s ^ {- 1}) d \mu (s).\tag{11}
\]

\begin{enumerate}
\def\labelenumi{\arabic{enumi})}
\setcounter{enumi}{1}
\tightlist
\item
  Let \(p\) and \(q\) be two conjugate exponents ( \(1 \leqslant p \leqslant +\infty\) ). If \(\mu \in \mathcal{M}^{\chi^{-1/q}}(\mathrm{G})\) and \(f \in \mathrm{L}^p(\mathrm{G},\beta)\) , then \(\mu\) and \(f\) are convolvable, and \(\mu * f\) is equal locally \(\beta\) -almost everywhere to a function in \(\mathrm{L}^p(\mathrm{G},\beta)\) ; one has (with an abuse of notations already noted)
\end{enumerate}

\[
\| \mu * f \| _ {p} \leqslant \| \mu \| _ {\chi^ {- 1 / q}} \| f \| _ {p}.
\]

If \(\mu \in \mathcal{M}^{\chi'^{-1/q}}(\mathrm{G})\) and \(f \in \mathrm{L}^p(\mathrm{G},\beta)\) , then \(f\) and \(\mu\) are convolvable, and \(f * \mu\) is equal locally \(\beta\) -almost everywhere to a function in \(\mathrm{L}^p(\mathrm{G},\beta)\) ; one has \(\| f * \mu \|_p \leqslant \| \mu \|_{\chi'^{-1/q}} \| f \|_p\) .

\begin{enumerate}
\def\labelenumi{\arabic{enumi})}
\setcounter{enumi}{2}
\item
  The mappings \((\mu, f) \mapsto \mu * f\) , \((f, \mu) \mapsto f * \mu\) define on \(\mathbf{L}^p(\mathbf{G}, \beta)\) the structures of a left module over \(\mathcal{M}^{\chi^{-1/q}}(\mathbf{G})\) and a right module over \(\mathcal{M}^{\chi'^{-1/q}}(\mathbf{G})\) . The two external laws on \(\mathbf{L}^p(\mathbf{G}, \beta)\) are permutable by the associativity of convolution.
\item
  If \(\mu * f\) is continuous and is given at every point by (10), then
\end{enumerate}

\[
(\mu * f) (e) = \int f (s ^ {- 1}) \chi (s ^ {- 1}) d \mu (s).\tag{12}
\]

If \(f * \mu\) is continuous and is given at every point by (11), then

\[
(f * \mu) (e) = \int f (s ^ {- 1}) \chi^ {\prime} (s ^ {- 1}) d \mu (s).\tag{13}
\]

\subsection{Convolution of functions on a group}

We conserve the notations G, \(\beta\) , \(\chi\) , \(\chi'\) of No.~4.

Recall that if f is a complex function on G, the property of being locally \(\beta\) -integrable is independent of the choice of \(\beta\) . Let \(\mathcal{L}(\mathrm{G})\) be the set of functions having this property. If \(f \in \mathcal{L}(\mathrm{G})\) , \(g \in \mathcal{L}(\mathrm{G})\) , the relation

\[\text{«}f \cdot \beta\text{ and }g \cdot \beta\text{ are convolvable »}\]

is independent of the choice of \(\beta\) (§3, No.~1, Prop. 6). We shall then say that \(f\) and \(g\) are convolvable. By No.~1, \((f\cdot \beta)*(g\cdot \beta)\) is of the form \(h\cdot \beta\) with \(h\in \mathcal{L}(\mathbf{G})\) , \(h\) being determined up to locally \(\beta\) -negligible sets. We shall write \(h = f * ^ {\beta} g\) and we shall say that \(h\) is a convolution product of \(f\) and \(g\) relative to \(\beta\) . (One omits \(\beta\) when no confusion is possible.) If \(\beta\) is replaced by \(\psi \cdot \beta\) , \(\psi\) being a continuous representation of \(\mathbf{G}\) in \(\mathbf{R}_+^*\) , \(h\) does not change (§3, No.~1, Prop. 6); if \(\beta\) is replaced by \(a\beta\) ( \(a\in \mathbf{R}_+^*\) ), then \(h\) is replaced by ah. The convolution product of several functions on G is defined in an analogous manner.

If one of the convolutions of f and g is continuous, it is uniquely determined since the support of \(\beta\) is G. It is then called the convolution product of f and g relative to \(\beta\) .

It is clear that

\[
f * ^ {\beta} g = (f \cdot \beta) * ^ {\beta} g = f * ^ {\beta} (g \cdot \beta).\tag{14}
\]

PROPOSITION 9. --- Let \(f,g\) be in \(\mathcal{L}(\mathbf{G})\) . Assume that the function \(s\mapsto g(s^{-1}x)f(s)\chi(s^{-1})\) is essentially \(\beta\) -integrable except for a locally \(\beta\) -negligible set of values of \(x\) , and that the function

\[
x \mapsto \int | g (s ^ {- 1} x) f (s) | \chi (s ^ {- 1}) d \beta (s),
\]

defined locally \(\beta\) -almost everywhere, is locally \(\beta\) -integrable. Then \(f\) and \(g\) are convolvable.

This follows from Prop. 1 of No.~1.

PROPOSITION 10. --- Let \(f, g\) be in \(\mathcal{L}(G)\) . Assume that one of these two functions is continuous or is zero on the complement of a countable union of compact sets. If \(f\) and \(g\) are convolvable, then the function \(f * g\) is given locally \(\beta\) -almost everywhere by

\[
\begin{array}{c} (f * g) (x) = \int_ {\mathrm{G}} g (s ^ {- 1} x) f (s) \chi (s ^ {- 1}) d \beta (s) \\ = \int_ {\mathrm{G}} f (x s ^ {- 1}) g (s) \chi^ {\prime} (s ^ {- 1}) d \beta (s). \end{array}\tag{15}
\]

This follows from Prop. 2 of No.~1, and the remarks in No.~4.

In particular, if \(f * g\) is continuous and is given at every point by (15), then

\[
(f * g) (e) = \int g (s ^ {- 1}) f (s) \chi (s ^ {- 1}) d \beta (s) = \int f (s ^ {- 1}) g (s) \chi^ {\prime} (s ^ {- 1}) d \beta (s).\tag{16}
\]

Still more particularly, if \(\beta\) is a left and right Haar measure, and if \(f*g\) and \(g*f\) are continuous and are given at every point by (15) and the analogous formula for \(g*f\) , then

\[
(f * g) (e) = (g * f) (e) = \int f (s) g (s ^ {- 1}) d \beta (s).\tag{17}
\]

PROPOSITION 11. --- Let \(f,g\) be in \(\mathcal{L}(\mathrm{G})\) . Assume that one of the functions \(f,g\) is continuous, and that one of the functions \(f,g\) has compact support. Then \(f\) and \(g\) are convolvable. The formula (15) defines for all \(x\in\mathrm{G}\) a product \(f*g\) that is continuous. If \(f\in\mathcal{K}(\mathrm{G})\) and \(g\in\mathcal{K}(\mathrm{G})\) , then \(f*g\in\mathcal{K}(\mathrm{G})\) .

This follows from Props. 3 and 4 of No.~2.

PROPOSITION 12. --- Let \(p\) and \(q\) be two conjugate exponents \((1 \leqslant p \leqslant +\infty)\) . If \(f\chi^{-1/q} \in \mathrm{L}^1(\mathrm{G}, \beta)\) and \(g \in \mathrm{L}^p(\mathrm{G}, \beta)\) , then \(f\) and \(g\) are convolvable, \(f * g\) is equal locally \(\beta\) -almost everywhere to a function in \(\mathrm{L}^p(\mathrm{G}, \beta)\) , and

\[
\left\| f * g \right\| _ {p} \leqslant \left\| f \chi^ {- 1 / q} \right\| _ {1} \left\| g \right\| _ {p}.
\]

If \(f \in \mathrm{L}^p(\mathbf{G},\beta)\) and \(g\chi'^{-1/q} \in \mathrm{L}^1(\mathbf{G},\beta)\) , then \(f\) and \(g\) are convolvable, \(f * g\) is equal locally \(\beta\) -almost everywhere to a function in \(\mathrm{L}^p(\mathbf{G},\beta)\) , and

\[
\| f * g \| _ {p} \leqslant \| f \| _ {p} \| g \chi^ {\prime - 1 / q} \| _ {1}.
\]

This follows from Props. 5 and 6 of No.~2 and the remarks in No.~4.

PROPOSITION 13. --- If \(f\chi^{-1} \in \mathbf{L}^1(\mathbf{G}, \beta)\) and \(g \in \overline{\mathcal{H}(\mathbf{G})}\) , or if \(f \in \overline{\mathcal{H}(\mathbf{G})}\) and \(g\chi'^{-1} \in \mathbf{L}^1(\mathbf{G}, \beta)\) , then \(f\) and \(g\) are convolvable, and (15) defines for every \(x \in \mathbf{G}\) a product \(f * g\) that belongs to \(\overline{\mathcal{H}(\mathbf{G})}\) .

This follows from Prop. 5 of No.~2, and the remarks in No.~4.

PROPOSITION 14. --- If \(f\chi^{-1} \in \mathbf{L}^1(\mathbf{G}, \beta)\) and \(g \in \mathbf{L}^\infty(\mathbf{G}, \beta)\) , then the formula (15) defines for every \(x \in \mathbf{G}\) a product \(f * g\) that is bounded and is uniformly continuous for the right uniform structure of \(\mathbf{G}\) .

We already know that \(f * g\) belongs to \(\mathrm{L}^{\infty}(\mathrm{G},\beta)\) (No.~2, Prop. 5); moreover, \((f*g)(x) = \int f(xs^{-1})g(s)d\nu(s)\) , on setting \(\nu = \chi'^{-1}\cdot \beta\) ; \(\nu\) is a right Haar measure. Therefore

\[
\begin{array}{c} | (f * g) (x) - (f * g) (x ^ {\prime}) | \leqslant \| g \| _ {\infty} \int | f (x s ^ {- 1}) - f (x ^ {\prime} s ^ {- 1}) | d \nu (s) \\ = \| g \| _ {\infty} \int | (f (s ^ {- 1}) - f (x ^ {\prime} x ^ {- 1} s ^ {- 1}) d \nu (s) \end{array}
\]

and the latter integral is arbitrarily small provided \(x'x^{-1}\) is in a suitable neighborhood of \(e\) (§2, No.~5, Prop. 8).

PROPOSITION 15. --- Let \(p\) and \(q\) be two conjugate exponents \((1 < p < +\infty)\) . Assume that \(\beta\) is left-invariant. Let \(f \in \mathrm{L}^p(\mathrm{G}, \beta)\) , \(g \in \mathrm{L}^q(\mathrm{G}, \breve{\beta})\) . Then \(f\) and \(g\) are convolvable. The formula (15) defines, for every \(x \in \mathrm{G}\) , a product \(f * g\) that belongs to \(\overline{\mathcal{H}(\mathrm{G})}\) and is such that \(\| f * g \|_{\infty} \leqslant \| f \|_p \| \breve{g} \|_q\) .

For, we have \(\check{g} \in \mathbf{L}^q(\mathbf{G}, \beta)\) , therefore the function \(s \mapsto g(s^{-1}x)f(s)\) is \(\beta\) -integrable for every \(x \in \mathbf{G}\) . Moreover,

\[
\begin{array}{c} \int | g (s ^ {- 1} x) f (s) | d \beta (s) \leqslant \bigg (\int | f (s) | ^ {p} d \beta (s) \bigg) ^ {1 / p} \bigg (\int | g (s ^ {- 1} x) | ^ {q} d \beta (s) \bigg) ^ {1 / q} \\ = \| f \| _ {p} \bigg (\int | \check {g} (x ^ {- 1} s) | ^ {q} d \beta (s) \bigg) ^ {1 / q} = \| f \| _ {p} \| \check {g} \| _ {q}, \end{array}
\]

therefore \(f\) and \(g\) are convolvable (Prop. 9). One sees at the same time that (15) defines for every \(x\) a product \(f * g\) such that

\[
\left| (f * g) (x) \right| \leqslant \| f \| _ {p} \| \check {g} \| _ {q}.
\]

For \(f, g\) in \(\mathcal{H}(\mathrm{G})\) , we have \(f * g \in \mathcal{H}(\mathrm{G})\) (Prop. 11); therefore, for \(f \in \mathrm{L}^p(\mathrm{G}, \beta)\) and \(g \in \mathrm{L}^q(\mathrm{G}, \beta)\) , the product \(f * g\) furnished by (15) is the uniform limit of functions in \(\mathcal{H}(\mathrm{G})\) , hence belongs to \(\overline{\mathcal{H}(\mathrm{G})}\) .

COROLLARY. --- Let \(f \in \mathbf{L}^2(\mathbf{G}, \beta)\) , \(g \in \mathbf{L}^2(\mathbf{G}, \beta)\) . Then \(f\) and \(\breve{g}\) are convolvable. One of the convolutions \(f * \breve{g}\) belongs to \(\overline{\mathcal{H}(\mathbf{G})}\) and its value at \(e\) is \(\int_{\mathbf{G}} f(s) \overline{g(s)} d\beta(s)\) .

It suffices to take \(p = q = 2\) in Prop. 15 and to apply (16).

We no longer assume \(\beta\) to be left-invariant. Let \(\rho\) be a lower semi-continuous finite function \textgreater0 on G, such that \(\rho(st) \leqslant \rho(s)\rho(t)\) for all s, t in G. We denote by \(\mathrm{L}^{\rho}(\mathrm{G}, \beta)\) the set of equivalence classes of the complex functions on G that are integrable for \(\rho \cdot \beta\) . By the mapping \(f \mapsto f \cdot \beta\) , \(\mathrm{L}^{\rho}(\mathrm{G}, \beta)\) may be identified with the set of elements of \(\mathcal{M}^{\rho}(\mathrm{G})\) that have base \(\beta\) (a set that is independent of the choice of \(\beta\) ). If one sets

\[
\| f \| _ {\rho} = \int_ {G} | f (s) | \rho (s) d \beta (s)
\]

for \(f \in \mathrm{L}^{\rho}(\mathrm{G},\beta)\) , this identification is compatible with the norms, thus \(\mathrm{L}^{\rho}(\mathrm{G},\beta)\) appears as a complete normed subalgebra of \(\mathcal{M}^{\rho}(\mathrm{G})\) . It is even a two-sided ideal of \(\mathcal{M}^{\rho}(\mathrm{G})\) by Prop. 10 of §3, No.~2. (For \(\rho = 1\) , one recovers one of the assertions of No.~4.) In particular, \(\mathrm{L}^{1}(\mathrm{G},\beta)\) may be identified with a closed two-sided ideal of \(\mathcal{M}^{1}(\mathrm{G})\) .

PROPOSITION 16. --- Let U be a continuous representation of G in a Banach space E. Set \(\rho(s) = \|U(s)\|\) for all \(s \in G\) . For every \(f \in L^{\rho}(G, \beta)\) , set \(U(f) = U(f \cdot \beta)\) . Then \(f \mapsto U(f)\) is a linear representation of the algebra \(L^{\rho}(G, \beta)\) in E, such that \(\|U(f)\| \leqslant \|f\|_{\rho}\) .

This follows from §2, No.~6 and §3, No.~3, Prop. 11.

\subsection{Applications}

PROPOSITION 17. --- Let G be a locally compact group, A a subset of G, measurable and not locally negligible for a Haar measure. Then \(A \cdot A^{-1}\) is a neighborhood of e.

Let \(\beta\) be a left Haar measure. There exists a compact subset \(K\) of \(G\) such that \(B = A \cap K\) is integrable with measure \(>0\) for \(\beta\) . Let us apply the Cor. of Prop. 15 with \(f = g = \varphi_{B}\) . The function \(F = \varphi_{B} * \check{\varphi}_{B}\) is continuous and \(>0\) at \(e\) . Therefore there exists a neighborhood \(V\) of \(e\) such that \(F(x) > 0\) for \(x \in V\) . Now,

\[
\mathrm{F} (x) = \int \varphi_ {\mathrm{B}} (s) \varphi_ {\mathrm{B}} (x ^ {- 1} s) d \beta (s) = \beta (\mathrm{B} \cap x \mathrm{B}).
\]

Therefore, for \(x \in V\) , one has \(B \cap xB \neq \varnothing\) , whence \(x \in B \cdot B^{-1}\) . Thus \(V \subset B \cdot B^{-1} \subset A \cdot A^{-1}\) .

COROLLARY 1. --- Let H be a subgroup of G measurable for a Haar measure \(\beta\) . Then H is either open or locally \(\beta\) -negligible.

For, \(H = H \cdot H^{-1}\) , therefore if H is not locally \(\beta\) -negligible, then H contains a neighborhood of e (Prop. 17) hence is open (GT, III, §2, No.~1, Cor. of Prop. 4).

COROLLARY 2. --- Let L be a subset of G stable for multiplication and whose complement is locally negligible for a Haar measure \(\beta\) . Then \(L = G\) .

For, \(L^{-1}\) and \(L \cap L^{-1}\) have locally \(\beta\) -negligible complements. Now, \(L \cap L^{-1}\) is a subgroup, hence is open (Cor. 1) and therefore closed. Therefore \(G - (L \cap L^{-1})\) , which is open and locally \(\beta\) -negligible, is empty. Thus \(G = L \cap L^{-1}\) .

PROPOSITION 18. --- Let G be a locally compact group, \(\Gamma\) a set equipped with a multiplication \((u,v)\mapsto uv\) and a Hausdorff topology such that:

\begin{enumerate}
\def\labelenumi{\arabic{enumi})}
\item
  the topology of \(\Gamma\) is invariant under the translations;
\item
  the restriction of the multiplication to every compact subset of \(\Gamma \times \Gamma\) is continuous.
\end{enumerate}

Let \(f: \mathbf{G} \to \Gamma\) be a mapping of \(\mathbf{G}\) into \(\Gamma\) such that \(f(xy) = f(x)f(y)\) for \(x, y\) in \(\mathbf{G}\) , and measurable for a Haar measure \(\beta\) on \(\mathbf{G}\) . Then \(f\) is continuous.

Set \(g(x) = f(x^{-1})\) for \(x \in \mathbf{G}\) . Since \(f\) and \(g\) are \(\beta\) -measurable, there exists a non \(\beta\) -negligible compact subset \(\mathbf{K}\) of \(\mathbf{G}\) such that the restrictions of \(f\) and \(g\) to \(\mathbf{K}\) are continuous. The mapping \((x,y) \mapsto f(xy^{-1}) = f(x)g(y)\) of \(\mathbf{K} \times \mathbf{K}\) into \(\Gamma\) is continuous because the multiplication of \(\Gamma\) is continuous on \(f(\mathbf{K}) \times g(\mathbf{K})\) ; now, this mapping may be written as \(\varphi \circ \psi\) , where \(\psi\) is the mapping \((x,y) \mapsto xy^{-1}\) of \(\mathbf{K} \times \mathbf{K}\) onto \(\mathbf{K} \cdot \mathbf{K}^{-1}\) , and \(\varphi\) is the restriction of \(f\) to \(\mathbf{K} \cdot \mathbf{K}^{-1}\) . Let \(\mathbf{R}\) be the equivalence relation defined on \(\mathbf{K} \times \mathbf{K}\) by \(\psi\) . The mapping \(\psi'\) of \((\mathbf{K} \times \mathbf{K}) / \mathbf{R}\) onto \(\mathbf{K} \cdot \mathbf{K}^{-1}\) deduced from \(\psi\) by passage to the quotient is continuous, therefore \((\mathbf{K} \times \mathbf{K}) / \mathbf{R}\) is Hausdorff and \(\psi'\) is a homeomorphism. Since \(\varphi \circ \psi\) is continuous, one sees that the restriction of \(f\) to \(\mathbf{K} \cdot \mathbf{K}^{-1}\) is continuous. Now, \(\mathbf{K} \cdot \mathbf{K}^{-1}\) is a neighborhood of \(e\) (Prop. 17), therefore \(f\) is continuous at \(e\) . For every \(x_0 \in \mathbf{G}\) , \(f(x_0x) = f(x_0)f(x)\) , thus \(f\) is continuous at \(x_0\) because the topology of \(\Gamma\) is invariant under translations.

COROLLARY 1. --- Let G be a locally compact group, \(\beta\) a Haar measure on G, E a Hausdorff barreled locally convex space, \(U\) a linear representation of G in E, such that \(U(s) \in \mathcal{L}(\mathrm{E};\mathrm{E})\) for all \(s \in \mathrm{G}\) , \(\beta\) -measurable when \(\mathcal{L}(\mathrm{E};\mathrm{E})\) is equipped with the topology of pointwise convergence. Then \(U\) is a continuous linear representation.

Let \(\Gamma\) be the group of automorphisms of E, equipped with the topology of pointwise convergence. This topology is Hausdorff and is invariant under translations. Let K be a compact subset of \(\Gamma\) . Then K is bounded in \(\mathcal{L}(\mathrm{E};\mathrm{E})\) equipped with the topology of pointwise convergence, hence is equicontinuous (TVS, III, §4, No.~2, Th. 1); therefore the mapping \((u,v)\mapsto v\circ u\) of \(\mathbf{K}\times \mathbf{K}\) into \(\mathcal{L}(\mathbf{E};\mathbf{E})\) is continuous (loc. cit., §5, No.~5, Cor. 1 of Prop. 9). Therefore, for every \(x\in \mathbf{E}\) , the mapping \(s\mapsto U(s)x\) of G into E is continuous (Prop. 18). Since E is barreled, \(U\) is continuous (§2, No.~1, Prop. 1).

COROLLARY 2. --- Let G be a locally compact group, \(\beta\) a Haar measure on G, E a separable Banach space, and \(U\) a linear representation of G in E such that \(U(s) \in \mathcal{L}(\mathrm{E};\mathrm{E})\) for all \(s \in \mathrm{G}\) . Let \((a_m)\) be a total sequence in E, and let \((a_n')\) be a dense sequence in the unit ball B' of the dual E' of E, equipped with the weak topology. Assume that the functions \(s \mapsto \langle U(s)a_m, a_n' \rangle\) on G are \(\beta\) -measurable. Then \(U\) is a continuous linear representation.

Let us first show that for every \(z' \in E'\) , the scalar functions

\[
s \mapsto \langle U (s) a _ {m}, z ^ {\prime} \rangle
\]

are \(\beta\) -measurable; we may restrict ourselves to the case that \(\|z'\| \leqslant 1\) , and, since \(B'\) is metrizable for the weak topology (TVS, III, §3, No.~4, Cor. 2 of Prop. 6), there exists a subsequence \((a_{n_k}')\) of \((a_n')\) that converges weakly to \(z'\) ; the function

\[
s \mapsto \langle U (s) a _ {m}, z ^ {\prime} \rangle
\]

is thus the limit of a sequence of \(\beta\) -measurable functions, whence our assertion. It follows that the mapping \(s \mapsto U(s)a_{m}\) of G into E is \(\beta\) -measurable for every m (Ch. IV, §5, No.~5, Prop. 10). On the other hand, there exists a sequence \((b_{m})\) of elements of E, linear combinations of the \(a_{i}\) , that is dense in the unit ball of E. For every \(s \in \mathbf{G}\) , \(\|U(s)\| = \sup_{m} \|U(s)b_{m}\|\) , therefore \(s \mapsto \|U(s)\|\) is measurable. Let K be a compact subset of G and let \(\varepsilon > 0\) . There exists a compact subset \(\mathrm{K}_{0}\) of K such that \(\beta(\mathrm{K} - \mathrm{K}_{0}) \leqslant \varepsilon\) and such that the restrictions to \(\mathrm{K}_{0}\) of the functions \(s \mapsto U(s)a_{m}\) and \(s \mapsto \|U(s)\|\) are continuous. Then the \(U(s)\) for \(s \in \mathrm{K}_{0}\) are equicontinuous, and the topology of pointwise convergence induces on \(U(\mathrm{K}_{0})\) the topology of pointwise convergence in the set of \(a_{m}\) (TVS, III, §3, No.~4, Prop. 5). Consequently the mapping \(s \mapsto U(s)\) of \(\mathrm{K}_{0}\) into \(\mathcal{L}_{s}(\mathrm{E};\mathrm{E})\) is continuous. It then suffices to apply Cor. 1.

\subsection{Regularization}

PROPOSITION 19. --- Let G be a locally compact group, \(\beta\) a relatively invariant positive measure \(\neq 0\) on G, B a base for the filter of neighborhoods of \(e\) in G, consisting of compact neighborhoods. For every V \(\in\) B, let \(f_{V}\) be a continuous function \(\geqslant 0\) on G, with support contained in V, such that \(\int f_{V} d\beta = 1\) . If \(\mu\) is a measure on G then, in \(\mathcal{M}(G)\) equipped with the topology of compact convergence in \(\mathcal{K}(G)\) ,

\[
\mu = \lim _ {\mathrm{V}} (\mu * f _ {\mathrm{V}}) \cdot \beta = \lim _ {\mathrm{V}} (f _ {\mathrm{V}} * \mu) \cdot \beta ,
\]

the limit being taken with respect to the section filter of B.

For the topology of compact convergence in \(\mathcal{C}(\mathrm{G})\) , \(f_{V} \cdot \beta\) tends to \(\varepsilon_{e}\) with respect to the section filter of B ( \(\S2\) , No.~7, Cor. 1 of Lemma 4). Therefore \(\mu = \lim_{\mathrm{V}} \mu * (f_{\mathrm{V}} \cdot \beta) = \lim_{\mathrm{V}} (f_{\mathrm{V}} \cdot \beta) * \mu\) in \(\mathcal{M}(\mathrm{G})\) equipped with the topology of compact convergence in \(\mathcal{K}(\mathrm{G})\) ( \(\S3\) , No.~3, Cor. of Prop. 12).

Remarks. --- 1) We thus see that every measure on G is the limit of measures admitting a continuous density with respect to every Haar measure (for the topology indicated in Prop. 19 and a fortiori for the vague topology).

\begin{enumerate}
\def\labelenumi{\arabic{enumi})}
\setcounter{enumi}{1}
\tightlist
\item
  If G is metrizable, B can be taken to be a sequence \((\mathrm{V}_{n})\) of neighborhoods. Then \(\mu\) is the limit of the sequence of measures \((\mu * f_{\mathrm{V}_{n}}) \cdot \beta\) with continuous densities. *If G is a real Lie group, the \(f_{V_{n}}\) can be taken to be infinitely differentiable; we shall see later on that the densities \(\mu * f_{V_{n}}\) are then infinitely differentiable.*
\end{enumerate}

PROPOSITION 20. --- We conserve the hypotheses and notations of Prop. 19. Let \(p \in [1, +\infty[\) and \(g \in \mathrm{L}^p(\mathbf{G}, \beta)\) . Then

\[
g = \lim _ {\mathrm{V}} g * ^ {\beta} f _ {\mathrm{V}} = \lim _ {\mathrm{V}} f _ {\mathrm{V}} * ^ {\beta} g
\]

in the sense of the norm \(N_{p}\) , the limit being taken with respect to the section filter of B.

It suffices to apply Prop. 6 (iii), and §2, No.~7, Cor. 3 of Lemma 4.

Remark 3. --- By Prop. 15, the functions \(g * f_{\mathrm{V}}\) , \(f_{\mathrm{V}} * g\) belong to \(\overline{\mathcal{K}(\mathrm{G})}\)

COROLLARY. --- Let W be a closed linear subspace of \(\mathrm{L}^{1}(\mathrm{G},\beta)\) . For W to be a left (resp. right) ideal of \(\mathrm{L}^{1}(\mathrm{G},\beta)\) , it is necessary and sufficient that W be invariant under the left (resp. right) translations of G.

Suppose that W is a left ideal. Let \(s \in G\) and \(g \in W\) . We have \(\varepsilon_{s} * g = \lim_{\mathbb{V}} f_{\mathbb{V}} * (\varepsilon_{s} * g) = \lim_{\mathbb{V}} (f_{\mathbb{V}} * \varepsilon_{s}) * g\) , and \((f_{\mathbb{V}} * \varepsilon_{s}) * g \in \mathbb{W}\) , therefore \(\varepsilon_{s} * g \in W\) , thus \(\boldsymbol{\gamma}(s)g \in \mathbb{W}\) . Conversely, if W is invariant under the left translations, then \(\mu *^{\beta}g \in W\) for \(\mu \in \mathcal{M}^{1}(\mathbf{G})\) and \(g \in W\) , therefore W is a fortiori a left ideal of \(\mathrm{L}^{1}(\mathrm{G}, \beta)\) . One argues similarly for right ideals.

Example. --- We take \(\mathbf{G} = \mathbf{R}\) . Let us define a function \(\mathbf{F}_n \in \mathcal{K}(\mathbf{R})\) by

\[
\begin{array}{l l} F _ {n} (x) = (1 - x ^ {2}) ^ {n} & \text { if } x \in [ - 1, 1 ] \\ F _ {n} (x) = 0 & \text { if } x \notin [ - 1, 1 ]. \end{array}
\]

Let \(\mathbf{A}_n = \int_{-1}^{+1}\mathbf{F}_n(x)dx\) , and \(\mathbf{G}_n = \mathbf{A}_n^{-1}\mathbf{F}_n\) . It is immediate that the measures \(\mathbf{G}_n(x)dx\) satisfy the conditions of §2, No.~7, Cor. 1 of Lemma 4. Let \(\mu\) be a measure on \(\mathbf{R}\) whose support is contained in \([-1/2, 1/2]\) . Then

\[
\begin{array}{r l} & {(\mu * \mathrm{G} _ {n}) (x) = \int_ {\mathbf {R}} \mathrm{G} _ {n} (x - y) d \mu (y)} \\ & {\qquad = \mathrm{A} _ {n} ^ {- 1} \int_ {- 1 / 2} ^ {1 / 2} \mathrm{F} _ {n} (x - y) d \mu (y).} \end{array}
\]

If \(-1 / 2\leqslant x\leqslant 1 / 2\) , then

\[
(\mu * \mathrm{G} _ {n}) (x) = \mathrm{A} _ {n} ^ {- 1} \int_ {- 1 / 2} ^ {1 / 2} [ 1 - (x - y) ^ {2} ] ^ {n} d \mu (y),
\]

therefore \(\mu * \mathbf{G}_n\) coincides on \([-1/2, 1/2]\) with a polynomial. In particular, if \(f\) is a continuous function with support contained in \([-1/2, 1/2]\) , then \(f * \mathbf{G}_n\) coincides in \([-1/2, 1/2]\) with a polynomial; moreover, by Prop. 5 (iv), and §2, No.~7, Cor. 3 of Lemma 4, \(f * \mathbf{G}_n\) converges uniformly to \(f\) . *If \(f\) is of class \(\mathbf{C}^r\) , the derivatives \(\mathrm{D}^s(f * \mathbf{G}_n)\) tend uniformly to \(\mathrm{D}^s f\) for \(0 \leqslant s \leqslant r\) .*

\section{THE SPACE OF CLOSED SUBGROUPS}

Throughout this section, G denotes a locally compact group and \(\mu\) a right Haar measure on G.

\subsection{The space of Haar measures on the closed subgroups of G}

Lemma 1. --- Let \(\alpha\) be a positive measure \(\neq 0\) on G, S its support; the following two conditions are equivalent:

\begin{enumerate}
\def\labelenumi{\alph{enumi})}
\item
  S is a closed subgroup of G and the measure induced by \(\alpha\) on S is a right Haar measure on S.
\item
  \(\delta(s)\alpha = \alpha\) for every \(s \in S\) .
\end{enumerate}

Moreover, when these conditions are satisfied, the set of \(t \in \mathbf{G}\) such that \(\delta(t)\alpha = \alpha\) is equal to \(S\) .

It is clear that a) implies b); conversely, the relation b) implies that \(Sx = S\) for every \(x \in S\) ; in other words, the relations \(x \in S\) and \(y \in S\) imply that \(y \in Sx\) , or again that \(yx^{-1} \in S\) , and since \(S\) is nonempty, \(S\) is a closed subgroup of \(G\) . The set of \(t \in G\) such that \(St = S\) is then equal to \(S\) itself, whence the last assertion.

For the rest of the section, we denote by \(\Gamma\) the set of positive measures \(\neq0\) on G satisfying the conditions of Lemma 1, and for every \(\alpha\in\Gamma\) we denote by \(H_{\alpha}\) the closed subgroup of G that is the support of \(\alpha\) .

PROPOSITION 1. --- The set \(\Gamma\) is closed in the space \(\mathcal{M}_{+}(G) = \{0\}\) equipped with the vague topology.

We first prove the following lemmas:

Lemma 2. --- Let X be a locally compact space and for every measure \(\alpha \in \mathcal{M}_{+}(X) - \{0\}\) , let \(S_{\alpha}\) be the support of \(\alpha\) . Let \(\Phi\) be a filter on \(\mathcal{M}_{+}(X) - \{0\}\) that converges vaguely to a measure \(\alpha_{0} \neq 0\) . Then, for every neighborhood V of a point s of the support of \(\alpha_{0}\) , there exists a set \(M \in \Phi\) such that, for every \(\alpha \in M\) , one has \(V \cap S_{\alpha} \neq \varnothing\) .

For, if \(\varphi \in \mathcal{K}_{+}(\mathbf{X})\) is a function with support contained in \(\mathbf{V}\) and such that \(\int \varphi(x)d\alpha_0(x) > 0\) , by definition there exists a set \(\mathbf{M} \in \Phi\) such that \(\int \varphi(x)d\alpha(x) > 0\) for all \(\alpha \in \mathbf{M}\) , which implies \(\mathbf{V} \cap \mathbf{S}_{\alpha} \neq \emptyset\) .

Lemma 3. --- Let E be a set filtered by a filter \(\Phi\) , and let \(\xi \mapsto \alpha(\xi)\) be a mapping of E into \(\Gamma\) that converges vaguely with respect to \(\Phi\) to a measure \(\alpha_0 \neq 0\) . On the other hand, let \(\xi \mapsto t_{\xi}\) be a mapping of E into G such that \(t_{\xi} \in \mathrm{H}_{\alpha(\xi)}\) for every \(\xi \in \mathrm{E}\) . If \(s\) is a cluster point of the mapping \(\xi \mapsto t_{\xi}\) with respect to \(\Phi\) , then \(\delta(s)\alpha_0 = \alpha_0\) .

Replacing if necessary \(\Phi\) by a finer filter, we can suppose that \(s\) is a limit of \(\xi \mapsto t_{\xi}\) with respect to \(\Phi\) ; by Lemma 1, \(\delta(t_{\xi})\alpha(\xi) = \alpha(\xi)\) for every \(\xi \in E\) , and the conclusion follows from the continuity of the mapping \((u,\lambda) \mapsto \delta(u)\lambda\) on \(G \times \mathcal{M}_{+}(G)\) (§3, No.~3, Prop. 13).

To prove Prop. 1 it suffices, by Lemma 1, to show that if a filter \(\Psi\) on \(\Gamma\) converges vaguely to a measure \(\alpha_0 \neq 0\) and if \(s\) belongs to the support of \(\alpha_0\) , then \(\delta(s)\alpha_0 = \alpha_0\) . Now, for every neighborhood \(V\) of \(s\) in \(G\) , there exists an \(M \in \Psi\) such that, for every \(\alpha \in M\) , one has \(V \cap H_\alpha \neq \varnothing\) , by Lemma 2. For every neighborhood \(V\) of \(s\) and every \(\alpha \in \Gamma\) , let \(t_{V,\alpha}\) be a point of \(V \cap H_\alpha\) if \(V \cap H_\alpha \neq \varnothing\) , and any point of \(H_\alpha\) in the contrary case; if \(\Theta\) is the section filter of the filter of neighborhoods of \(s\) , and \(\Phi\) is the product filter \(\Theta \times \Psi\) , then \(s\) is, by the foregoing, a cluster point of \((V,\alpha) \mapsto t_{V,\alpha}\) with respect to \(\Phi\) . Since, on the other hand, the mapping \((V,\alpha) \mapsto \alpha\) has \(\alpha_0\) as limit with respect to \(\Phi\) , the proposition follows from Lemma 3.

PROPOSITION 2. --- Let \(\varphi\) be a function in \(\mathcal{K}_{+}(\mathbf{G})\) such that \(\varphi(e) > 0\) . Then the set \(\Gamma_{\varphi}\) of measures \(\alpha \in \Gamma\) such that \(\int \varphi(x) d\alpha(x) = 1\) is compact for the vague topology.

The set \(\Gamma_{\varphi}\) is the intersection of \(\Gamma\) with the hyperplane of \(\mathcal{M}(\mathrm{G})\) formed by the \(\alpha\) such that \(\int \varphi(x)d\alpha(x)=1\) ; since this hyperplane is vaguely closed in \(\mathcal{M}(\mathrm{G})\) and does not contain 0, it follows from Prop. 1 that \(\Gamma_{\varphi}\) is vaguely closed in \(\mathcal{M}(\mathrm{G})\) . It therefore suffices to show that for every compact subset K of G, one has \(\sup_{\alpha\in\Gamma_{\varphi}}\alpha(\mathrm{K})<+\infty\) (Ch. III, §1,

No.~9, Prop. 15). Now, let U be the open neighborhood of e in G defined by the inequality \(\varphi(x) > \varphi(e)/2\) ; since \(1 = \int \varphi(x) d\alpha(x) \geqslant \int_{\mathrm{U}} \varphi(x) d\alpha(x)\) for \(\alpha \in \Gamma_{\varphi}\) , one sees that, on setting \(c = 2/\varphi(e)\) , one has \(\alpha(\mathrm{U}) \leqslant c\) for every \(\alpha \in \Gamma_{\varphi}\) . Let V be a symmetric open neighborhood of e in G such that \(V^{2} \subset U\) ; let us show that \(\alpha(Vx) \leqslant c\) for every \(x \in G\) and every \(\alpha \in \Gamma_{\varphi}\) . Indeed, this relation is trivial if Vx does not intersect the support \(H_{\alpha}\) of \(\alpha\) ; if, on the contrary, there exists an \(h \in Vx \cap H_{\alpha}\) , then h = vx for some \(v \in V\) , whence

\[
\mathbf {V} x = \mathbf {V} v ^ {- 1} h \subset \mathbf {V} ^ {2} h \subset \mathrm{U} h,
\]

and since \(\delta(h)\alpha = \alpha\) , it follows that \(\alpha(\mathrm{V}x) \leqslant \alpha(\mathrm{U}h) = \alpha(\mathrm{U}) \leqslant c\) . Now let \((x_i)_{1 \leqslant i \leqslant n}\) be a sequence of points of K such that the \(Vx_i\) form a covering of K; it follows from the foregoing that \(\alpha(K) \leqslant \sum_{i=1}^{n} \alpha(Vx_i) \leqslant nc\) for every \(\alpha \in \Gamma_\varphi\) ; Q.E.D.

PROPOSITION 3. --- Under the hypotheses of Prop. 2, the mapping \(\alpha \mapsto \left( \langle \varphi, \alpha \rangle, \frac{\alpha}{\langle \varphi, \alpha \rangle} \right)\) is a homeomorphism of \(\Gamma\) onto the product space \(\mathbf{R}_{+}^{*} \times \Gamma_{\varphi}\) .

Since the mapping \(\alpha \mapsto \langle \varphi, \alpha \rangle\) is vaguely continuous, it suffices to observe that \(\langle \varphi, \alpha \rangle \neq 0\) for every measure \(\alpha \in \Gamma\) , since \(e\) belongs to the support \(\mathbf{H}_{\alpha}\) of \(\alpha\) and \(\varphi(e) > 0\) .

\subsection{Semi-continuity of the volume of the homogeneous space}

In this No., for every measure \(\alpha \in \Gamma\) we set

\[
\mathrm{Q} _ {\alpha} = \mathrm{G} / \mathrm{H} _ {\alpha},\tag{1}
\]

and we denote by \(\pi_{\alpha}\) the canonical mapping \(\mathbf{G} \to \mathbf{Q}_{\alpha}\) .

Let \(\Gamma^0\) be the subset of \(\Gamma\) formed by the measures \(\alpha\) such that the subgroup \(\mathbf{H}_{\alpha}\) of \(\mathbf{G}\) is unimodular; the elements of \(\Gamma^0\) are characterized by the fact that \(\alpha(f) = \alpha(\check{f})\) for every function \(f \in \mathcal{K}(\mathbf{G})\) (every function of \(\mathcal{K}(\mathbf{H}_{\alpha})\) being extendible to a function of \(\mathcal{K}(\mathbf{G})\) by Urysohn's theorem); it follows that \(\Gamma^0\) is a closed subset of \(\Gamma\) . Recall that for every \(\alpha \in \Gamma^0\) , the quotient measure \(\mu_{\alpha} = \mu / \alpha\) on \(\mathbf{Q}_{\alpha}\) is defined and is relatively invariant under \(\mathbf{G}\) (Ch. VII, §2, No.~6, Th. 3); also recall that for every function \(f \in \mathcal{K}(\mathbf{G})\) ,

\[
\int_ {\mathbf {G}} f (x) d \mu (x) = \int_ {\mathbf {Q} _ {\alpha}} d \mu_ {\alpha} (\dot {x}) \int_ {\mathrm{H} _ {\alpha}} f (x s) d \alpha (s),\tag{2}
\]

where \(\dot{x} = \pi_{\alpha}(x)\) is the canonical image of \(x\in \mathbf{G}\) in \(\mathbf{Q}_{\alpha}\) .

PROPOSITION 4. --- Let \(\Gamma^0\) be the set of measures \(\alpha \in \Gamma\) such that \(\mathrm{H}_{\alpha}\) is unimodular, and for every \(\alpha \in \Gamma^0\) set \(\mu_{\alpha} = \mu / \alpha\) ; then the mapping \(\alpha \mapsto \| \mu_{\alpha} \|\) of \(\Gamma^0\) into \(\overline{\mathbf{R}}\) is lower semi-continuous for the vague topology.

For every \(\alpha \in \Gamma^0\) and every function \(f\in \mathcal{K}(\mathbf{G})\) , set

\[
f _ {\alpha} (\dot {x}) = \int_ {\mathrm{H} _ {\alpha}} f (x s) d \alpha (s) = (f * \alpha) (x),
\]

where the convolution product is taken relative to the right Haar measure \(\mu\) and where one makes use of the fact that \(\check{\alpha} = \alpha\) (§4, No.~4, formula (11)). We know (Ch. VII, §2, No.~1, Prop. 2) that the mapping \(f \mapsto f_{\alpha}\) of \(\mathcal{H}_{+}(G)\) into \(\mathcal{H}_{+}(Q_{\alpha})\) is surjective; therefore, by (2),

\[
\| \mu_ {\alpha} \| = \sup _ {f \in \mathcal {K} _ {+} (\mathrm{G}), f \neq 0} \mu_ {\alpha} (f _ {\alpha}) / \| f _ {\alpha} \| = \sup _ {f \in \mathcal {K} _ {+} (\mathrm{G}), f \neq 0} \mu (f) / \| f _ {\alpha} \|,
\]

where one has set

\[
\left\| f _ {\alpha} \right\| = \sup _ {\dot {x} \in \mathrm{Q} _ {\alpha}} | f _ {\alpha} (\dot {x}) | = \sup _ {x \in \mathrm{G}} | (f * \alpha) (x) |.\tag{3}
\]

To establish the proposition, it will suffice to show that, given \(f \in \mathcal{K}_{+}(G)\) , the mapping \(\alpha \mapsto \|f_{\alpha}\|\) is vaguely continuous. Now, let K be the support of f; the function \(f * \alpha\) has its support contained in \(KH_{\alpha}\) and is invariant on the right under \(H_{\alpha}\) ; consequently

\[
\left\| f _ {\alpha} \right\| = \sup _ {x \in K} | (f * \alpha) (x) |.
\]

The conclusion therefore follows from the fact that the mapping \(\alpha\mapsto f*\alpha\) of \(\mathcal{M}_{+}(\mathrm{G})\) equipped with the vague topology, into \(\mathcal{C}(\mathrm{G})\) equipped with the topology of compact convergence, is continuous (§4, No.~2, Remark 1).

Recall that if, for a measure \(\alpha \in \Gamma^0\) , \(\| \mu_{\alpha}\|\) is finite, then G is necessarily unimodular (Ch. VII, §2, No.~6, Cor. 3 of Th. 3).

PROPOSITION 5. --- Let g be a \(\mu\) -integrable positive numerical function and let \(\Gamma^{0}(g)\) be the set of measures \(\alpha\in\Gamma^{0}\) such that \(\int^{*}g(xs)d\alpha(s)\geqslant1\) for all \(x\in G\) . Then the mapping \(\alpha\mapsto\|\mu_{\alpha}\|\) of \(\Gamma^{0}(g)\) into \(\overline{\mathbf{R}}\) is vaguely continuous.

For every measure \(\alpha \in \Gamma^0 (\mathbf{G})\) , recall (Ch. VII, §2, No.~3, Prop. 5) that the function

\[
g _ {\alpha} (\dot {x}) = \int_ {\mathrm{H} _ {\alpha}} g (x s) d \alpha (s)
\]

is defined \(\mu_{\alpha}\) -almost everywhere on \(Q_{\alpha}\) , is \(\mu_{\alpha}\) -integrable, and

\[
\int_ {G} g (x) d \mu (x) = \int_ {Q _ {\alpha}} g _ {\alpha} (\dot {x}) d \mu_ {\alpha} (\dot {x}).\tag{4}
\]

In view of Prop. 4, it suffices to prove that, in \(\Gamma^0(g)\) , \(\alpha \mapsto \| \mu_\alpha \|\) is upper semi-continuous. Fix a measure \(\alpha \in \Gamma^0(g)\) , and let \(K\) be a compact subset of \(G\) . There exists on \(Q_\alpha\) a continuous function with compact support, taking its values in {[}0,1{]}, equal to 1 on the compact set \(\pi_\alpha(K)\) ; since the mapping \(f \mapsto f_\alpha\) of \(\mathcal{H}_+(\mathrm{G})\) into \(\mathcal{H}_+(\mathrm{Q}_\alpha)\) is surjective (Ch. VII, §2, No.~1, Prop. 2), one sees that there exists a function \(f \in \mathcal{H}_+(\mathrm{G})\) such that

\[
(f * \alpha) (x) = \int_ {\mathrm{G}} f (x s)   d \alpha (s) \left\{ \begin{array}{l} \leqslant 1 \text {for all} x \in \mathrm{G} \\ = 1 \text {for all} x \in \mathrm{K}. \end{array} \right.
\]

Since \(\beta \mapsto f * \beta\) is a continuous mapping of \(\mathcal{M}_{+}(\mathbf{G})\) , equipped with the vague topology, into \(\mathcal{C}(\mathbf{G})\) equipped with the topology of compact convergence (§4, No.~2, Remark 1), one sees that for every \(\varepsilon > 0\) , the set \(\mathrm{U}_{\varepsilon}\) of \(\beta \in \Gamma^0 (\mathbf{G})\) such that

\[
f _ {\beta} (\dot {x}) = \int_ {\mathrm{G}} f (x s) d \beta (s) > 1 - \varepsilon \quad \text { for   all } x \in \mathrm{K}
\]

is an open neighborhood of \(\alpha\) in \(\Gamma^0 (g)\) ; for every \(\beta \in \mathrm{U}_{\varepsilon}\) , we then have, by virtue of the formula (2),

\[
\left\| \mu_ {\alpha} \right\| \geqslant \int_ {G} f (x) d \mu (x) = \int_ {Q _ {\beta}} f _ {\beta} (\dot {x}) d \mu_ {\beta} (\dot {x}) \geqslant (1 - \varepsilon) \mu_ {\beta} \left(\pi_ {\beta} (K)\right).\tag{5}
\]

Given a number \(\varepsilon > 0\) , let us choose a function \(h \in \mathcal{H}_{+}(\mathrm{G})\) such that \(\int_{\mathrm{G}} |g(x) - h(x)| d\mu(x) \leqslant \varepsilon\) , and let us take \(\mathrm{K} = \operatorname{Supp}(h)\) in the foregoing. For every \(\beta \in \Gamma^0(g)\) , by hypothesis \(g_\beta(\dot{x}) \geqslant 1\) almost everywhere (for \(\mu_\beta\) ) in \(\mathbf{Q}_\beta\) , therefore

\[
\mu_ {\beta} \big (\mathrm{Q} _ {\beta} - \pi_ {\beta} (\mathrm{K}) \big) \leqslant \int_ {\mathrm{Q} _ {\beta} - \pi_ {\beta} (\mathrm{K})} g _ {\beta} (\dot {x}) d \mu_ {\beta} (\dot {x}) = \int_ {\mathrm{G} - \mathrm{KH} _ {\beta}} g (x) d \mu (x)
\]

by virtue of (4); since \(h\) is zero outside \(\mathbf{K}\) , and a fortiori outside \(\mathrm{KH}_{\beta}\) , it follows that

\[
\begin{array}{c} \mu_ {\beta} \big (\mathrm{Q} _ {\beta} - \pi_ {\beta} (\mathrm{K}) \big) \leqslant \int_ {\mathrm{G} - \mathrm{KH} _ {\beta}} | g (x) - h (x) | d \mu (x) \\ \leqslant \int_ {\mathrm{G}} | g (x) - h (x) | d \mu (x) \leqslant \varepsilon ; \end{array}
\]

combining this result with (5), one sees that

\[
\| \mu_ {\beta} \| \leqslant \varepsilon + \| \mu_ {\alpha} \| / (1 - \varepsilon)
\]

when \(\beta \in \mathrm{U}_{\varepsilon}\) , which completes the proof.

COROLLARY 1. --- Let K be a compact subset of G, V a symmetric compact neighborhood of e in G, c a real number \textgreater{} 0. The restriction of the mapping \(\alpha \mapsto \|\mu_{\alpha}\|\) to the set of \(\alpha \in \Gamma^0\) such that \(G = KH_{\alpha}\) and \(\alpha(V) \geqslant c\) is vaguely continuous.

For, let \(g \in \mathcal{K}_{+}(\mathbf{G})\) be a function such that \(g(x) \geqslant 1 / c\) for \(x \in \mathrm{KV}\) . For every \(x \in \mathbf{K}\) ,

\[
\int g (x s) d \alpha (s) \geqslant \int_ {\mathrm{V}} g (x s) d \alpha (s) \geqslant 1
\]

for \(\alpha\) satisfying the conditions of the statement; since, moreover, \(\pi_{\alpha}(\mathbf{K}) = \mathrm{Q}_{\alpha}\) , one has \(\alpha \in \Gamma^0(g)\) , whence the corollary.

COROLLARY 2. --- Let A be a \(\mu\) -integrable subset of G. The restriction of the mapping \(\alpha \mapsto \|\mu_{\alpha}\|\) to the set \(N_A\) of normalized Haar measures of the discrete subgroups H of G such that \(G = AH\) , is vaguely continuous.

For \(a \in \mathbf{A}\) and \(\alpha \in \mathbf{N}_{\mathbf{A}}\) ,

\[
\int \varphi_ {\mathrm{A}} (a s) d \alpha (s) \geqslant \varphi_ {\mathrm{A}} (a) = 1,
\]

and since \(\pi_{\alpha}(\mathbf{A}) = \mathbf{Q}_{\alpha}\) , one has \(\mathrm{N}_{\mathrm{A}} \subset \Gamma^{0}(\varphi_{\mathrm{A}})\) , and the corollary therefore follows from Prop. 5.

\subsection{The space of closed subgroups of G}

Let us denote by \(\Sigma\) the set of closed subgroups of G; if one associates to each measure \(\alpha\in\Gamma\) the subgroup \(H_{\alpha}\) that is the support of \(\alpha\) , one obtains a mapping (called canonical) of \(\Gamma\) into \(\Sigma\) , which is clearly surjective and permits canonically identifying \(\Sigma\) with the set of orbits of the group of homotheties in \(\Gamma\) with ratio \textgreater0. The set \(\Sigma\) , equipped with the quotient topology of the vague topology on \(\Gamma\) , is called the space of closed subgroups of G.

THEOREM 1. --- Let G be a locally compact group. The space \(\Sigma\) of closed subgroups of G is compact. Moreover, one has the following properties:

\begin{enumerate}
\def\labelenumi{(\roman{enumi})}
\item
  The set \(\Sigma^0\) of unimodular closed subgroups of \(\mathbf{G}\) is closed in \(\Sigma\) (hence is compact).
\item
  If G is generated by a compact neighborhood of e, then the set \(\Sigma_{c}^{0}\) of unimodular closed subgroups H of G such that the quotient space G/H is compact, is open in \(\Sigma^{0}\) (hence is locally compact).
\item
  For every relatively compact open neighborhood U of e in G, the set \(D_{U}\) of discrete subgroups H of G such that \(H \cap U = \{e\}\) is closed in \(\Sigma^{0}\) (hence is compact).
\end{enumerate}

It follows from Prop. 3 of No.~1 that \(\Sigma\) is homeomorphic to \(\Gamma_{\varphi}\) , hence is compact by Prop. 2 of No.~1. Moreover, it was noted at the beginning of No.~2 that the set \(\Gamma^0\) of measures \(\alpha \in \Gamma\) such that \(\mathbf{H}_{\alpha}\) is unimodular is closed in \(\Gamma\) ; since \(\Gamma^0\) is stable under the homotheties with ratio \(>0\) , the image \(\Sigma^0\) of \(\Gamma^0\) in \(\Sigma\) is a closed subset of \(\Sigma\) , which proves (i).

Property (ii) will be a consequence of the following proposition:

PROPOSITION 6. --- Suppose that the locally compact group G is generated by a compact neighborhood of e. Then the set \(\Gamma_{c}^{0}\) of measures \(\alpha\in\Gamma^{0}\) such that \(\mathbf{G} / \mathbf{H}_{\alpha}\) is compact is open in \(\Gamma^0\) , and the restriction to \(\Gamma_c^0\) of the mapping \(\alpha \mapsto \| \mu_{\alpha}\|\) is vaguely continuous.

With the notations of Prop. 5 of No.~2, we have, for \(g \in \mathcal{K}_{+}(\mathbf{G})\) ,

\[
\Gamma^ {0} (g) \subset \Gamma_ {c} ^ {0}.\tag{6}
\]

For, if K is the support of g, the relation \(\int g(xs) d\alpha(s) \geqslant 1\) for all \(x \in G\) implies \(KH_{\alpha} = G\) , the integral obviously being zero on the complement of \(KH_{\alpha}\) , therefore \(G/H_{\alpha} = \pi_{\alpha}(K)\) is compact. Given a measure \(\alpha \in \Gamma_{c}^{0}\) , it will therefore suffice to define a function \(g \in \mathcal{K}_{+}(G)\) such that \(\Gamma^{0}(g)\) is a neighborhood of \(\alpha\) in \(\Gamma^{0}\) . Since \(G/H_{\alpha}\) is compact and the canonical mapping \(f \mapsto f_{\alpha}\) of \(\mathcal{K}_{+}(G)\) into \(\mathcal{K}_{+}(G/H_{\alpha})\) is surjective (Ch. VII, §2, No.~2), there exists a function \(g \in \mathcal{K}_{+}(G)\) such that \(\int g(xs) d\alpha(s) = 2\) for every \(x \in G\) . Let K be the (compact) support of g, L a symmetric compact neighborhood of e in G that generates G; the mapping \(\beta \mapsto g*\beta\) of \(\mathcal{M}_{+}(G)\) into \(\mathcal{C}(G)\) being vaguely continuous (§4, No.~2, Remark 1), there exists a neighborhood W of \(\alpha\) in \(\Gamma^{0}\) such that

\[
(g * \beta) (x) = \int g (x s) d \beta (s) \geqslant 1\tag{7}
\]

for all \(\beta \in W\) and \(x \in LK\) . The first member of (7) being equal to zero outside \(\mathrm{KH}_{\beta}\) , the relation \(\beta \in W\) implies

\[
\mathbf {L K} \subset \mathbf {K H} _ {\beta},
\]

from which one deduces, by induction on n, that \(L^{n}K \subset KH_{\beta}\) for every integer n \textgreater{} 0; since L generates G, we therefore have \(G = KH_{\beta}\) for every measure \(\beta \in W\) , which proves that \(W \subset \Gamma_{c}^{0}\) . On the other hand, the first member of (7) being invariant on the right under \(H_{\beta}\) , the inequality (7) is also valid for \(x \in LKH_{\beta} = G\) ; therefore \(W \subset \Gamma^{0}(g)\) , which proves the proposition.

Finally, (iii) will be a consequence of the following proposition:

PROPOSITION 7. --- Let \(N \subset \Gamma^{0}\) be the subspace of normalized Haar measures on the discrete subgroups of G, and for every relatively compact open neighborhood U of e in G, let \(N_{U}\) be the subset of N formed by the \(\alpha\) such that \(H_{\alpha} \cap U = \{e\}\) . Then:

\begin{enumerate}
\def\labelenumi{\alph{enumi})}
\item
  \(\mathbf{N}_{\mathrm{U}}\) is compact.
\item
  The interiors of the sets \(\mathbf{N}_{\mathbf{U}}\) in \(\mathbf{N}\) form a covering of \(\mathbf{N}\) , as \(\mathbf{U}\) runs over the set of relatively compact open neighborhoods of \(e\) in \(\mathbf{G}\) .
\item
  For a subset M of N to be relatively compact in N, it is necessary and sufficient that there exist a relatively compact open neighborhood U of e in G such that \(M \subset N_{U}\) .
\end{enumerate}

Since \(D_{U}\) is the image of \(N_{U}\) under the canonical continuous mapping \(\Gamma \rightarrow \Sigma\) , the assertion (iii) of Th. 1 will result at once from Prop. 7 a).

To prove Prop. 7, we observe that \(N_{U}\) can be defined as the subset of \(\Gamma^{0}\) formed by the \(\alpha\) such that both

\[
\alpha (\{e \}) \geqslant 1 \quad \text { and } \quad \alpha (U) \leqslant 1.
\]

Now, if A is compact (resp. open and relatively compact) in G, then the mapping \(\alpha \mapsto \alpha(A)\) of \(\mathcal{M}_{+}(G)\) into R is upper (resp. lower) semicontinuous for the vague topology (Ch. IV, §4, No.~4, Cor. 3 of Prop. 5 and loc. cit., §1, No.~1, Prop. 4); we thus see that \(N_{U}\) is a closed subset of \(\Gamma^{0}\) . Moreover, let \(\varphi \in \mathcal{K}_{+}(G)\) be a function such that \(\varphi(e) = 1\) and \(\varphi(x) = 0\) on G - U; it is clear that \(\int \varphi(x)d\alpha(x) = 1\) for all \(\alpha \in N_{U}\) ; Prop. 2 of No.~1 therefore shows that \(N_{U}\) is a compact set, which proves a). On the other hand let V be a relatively compact open neighborhood of \(e\) in G such that \(\overline{V} \subset U\) , and let \(\varphi \in \mathcal{K}_{+}(G)\) , with support contained in U and such that \(\varphi(x) = 1\) on V. Then \(\alpha(\varphi) = 1\) for \(\alpha \in N_{U}\) , therefore there exists a neighborhood W of \(\alpha\) in N such that \(\beta(\varphi) < 2\) for \(\beta \in W\) ; it is then clear that \(W \subset N_{V}\) , therefore \(N_{V}\) is a neighborhood of \(N_{U}\) . Since the \(N_{U}\) cover N, this proves b). Finally, every compact subset M of N is contained in a finite union of sets \(N_{U_i}\) ( \(1 \leqslant i \leqslant n\) ), and since \(\bigcup_{i} N_{U_i} \subset N_{U}\) , where \(U = \bigcap_{i} U_i\) , this proves c).

COROLLARY. --- The subspace N of \(\Gamma^0\) is locally compact.
